% \documentclass[acmsmall,screen,review,anonymous,nonacm]{acmart}
\documentclass[acmsmall,screen,nonacm]{acmart}

% --------------------------------------
% Recommended ACM packages
% --------------------------------------
\usepackage{booktabs}      % high-quality tables
\usepackage{listings}      % code listings
\usepackage{xcolor}
\usepackage{graphicx}
\usepackage{rotating} % optional, for \rotatebox if not already available
\usepackage{algorithm2e}
\usepackage{tikz}
\usetikzlibrary{arrows.meta,positioning,calc}
\usepackage{makecell}
\usepackage{booktabs}
\usepackage{caption}
\usepackage{ragged2e}
\usepackage{multirow}
\usepackage{array} % for @{} column spacing control
\usepackage{cleveref}
\usepackage{placeins}
\usepackage{minted}
%\usepackage[frozencache,cachedir=.]{minted} % for arxiv
%\usepackage[finalizecache,cachedir=.]{minted} % to generate cache
\usepackage{subcaption}
\usepackage[most]{tcolorbox}
\usepackage{svg}
\usepackage{wrapfig}
\usepackage{tcolorbox}
\usepackage{etoolbox}
\tcbuselibrary{breakable}
\usepackage{placeins} % in preamble

\usepackage{enumitem}
\setlist{leftmargin=7mm}
\usepackage{xcolor}

\newif\ifrevision
% \revisiontrue   % <-- turn ON revisions
\revisionfalse % <-- turn OFF revisions
\ifrevision
  \newcommand{\rev}[1]{\textcolor{red}{#1}}
\else
  \newcommand{\rev}[1]{#1}
\fi

\ifrevision
  \newcommand{\temp}[1]{\textcolor{orange}{#1}}
\else
  \newcommand{\temp}[1]{#1}
\fi
\usepackage{xcolor}

% \newif\ifrevision
% % \revisiontrue   % turn OFF with \revisionfalse
% \revisionfalse

\ifrevision
  \newenvironment{revblock}{\color{red}}{}
\else
  \newenvironment{revblock}{}{}
\fi
\newtcbox{\hly}{on line,
    arc=0pt,
    colframe=yellow!50, % Border color
    colback=yellow!50, % Background color
    boxrule=0.0pt, % Border thickness
    boxsep=0pt, % Remove padding inside the box
    left=1pt, right=1pt, top=2pt, bottom=2pt % Minimal margins
}

\newtcbox{\hlgr}{on line,
    arc=0pt,
    colframe=gray!30, % Border color
    colback=gray!30, % Background color
    boxrule=0.0pt, % Border thickness
    boxsep=0pt, % Remove padding inside the box
    left=1pt, right=1pt, top=2pt, bottom=2pt % Minimal margins
}


\newtcbox{\hlg}{on line,
    arc=0pt,
    colframe=green!50, % Border color
    colback=green!50, % Background color
    boxrule=0.0pt, % Border thickness
    boxsep=0pt, % Remove padding inside the box
    left=1pt, right=1pt, top=2pt, bottom=2pt % Minimal margins
}


\tcbset{
    mybox/.style={colback=white, width=\textwidth, boxrule=0.2mm,title=Prompt, sharp corners,left=1mm, right=1mm, bottom=1mm},
}
\tcbset{enhanced,attach boxed title to top left={xshift=2mm,yshift=-2mm},coltitle=black,boxed title style={size=small,colback=white, boxrule=0.2mm}}


% --------------------------------------
% Your custom macros (SAFE TO KEEP)
% --------------------------------------

\newcommand{\erase}{\mathsf{erase}}
\newcommand{\shape}{\mathrm{shape}}
\newcommand{\Cfg}{\mathcal{C}}
\newcommand{\Vars}{\mathcal{V}}
\newcommand{\Scalars}{\mathcal{S}}
\newcommand{\Ptrs}{\mathcal{P}}
\newcommand{\Mem}{\mathcal{M}}
\newcommand{\Val}{\mathcal{Val}}
\newcommand{\Fields}{\mathcal{F}}
\newcommand{\NULL}{\mathsf{NULL}}
\newcommand{\projName}{\textsc{Gordian}\xspace}
\lstdefinelanguage{text}{
  moredelim=[is][\ttfamily]{\%\%}{\%\%},
  basicstyle=\ttfamily,
}

% --------------------------------------
% C-style listings (ACM-compatible)
% --------------------------------------
\lstdefinestyle{CStyle}{
  language=C,
  basicstyle=\ttfamily\footnotesize,
  keywordstyle=\color{blue},
  commentstyle=\color{gray},
  stringstyle=\color{orange!60!black},
  numbers=left,
  numberstyle=\tiny\color{gray},
  stepnumber=1,
  frame=single,
  xleftmargin=10pt,
  xrightmargin=10pt,
  breaklines=true,
  tabsize=2
}

% Shortcut
\newcommand{\codelist}[2]{%
  \begin{lstlisting}[style=CStyle,caption={#1}]
#2
  \end{lstlisting}
}

\newcommand{\sergey}[1]{{\color{red}\textbf{(Sergey:} #1\textbf{)}}}
\newcommand{\dimitris}[1]{{\color{red}\textbf{(Dimitris:} #1\textbf{)}}}


% --------------------------------------
% Theorem environments (ACM-safe)
% --------------------------------------
\theoremstyle{definition}
\newtheorem{definition}{Definition}
\newtheorem{theorem}{Theorem}
\newtheorem{corollary}{Corollary}

% --------------------------------------
% Metadata (ANONYMOUS FOR REVIEW)
% --------------------------------------
\title[Defusing Logic Bombs in Symbolic Execution with LLM-Generated Ghost Code]{Defusing Logic Bombs in Symbolic Execution\\ with LLM-Generated Ghost Code}

% Authors omitted for double-anonymous review
% \author{...}

\author{Dimitrios Stamatios Bouras}
\email{2501112125@stu.pku.edu.cn}
\affiliation{%
  \institution{Key Lab of HCST (PKU), MoE, SCS, Peking University}
  \city{Beijing}
  \country{China}
}

\author{Sergey Mechtaev}
\authornote{Sergey Mechtaev is the corresponding author.}
\email{mechtaev@pku.edu.cn}
\affiliation{%
  \institution{Key Lab of HCST (PKU), MoE, SCS, Peking University}
  \city{Beijing}
  \country{China}}


\begin{document}


% 30.86 klee improve fdlib
% Across these benchmarks, \projName{} improves the main coverage metric by
% 28.5--115.2\% over the strongest traditional symbolic-execution baseline
% and by 74.1--189.8\% over the strongest relevant LLM-based baseline.
\begin{abstract}
    Symbolic execution is a powerful program analysis technique, but its effectiveness is fundamentally limited by solver-hostile program fragments, complex numerical reasoning, and unbounded heap structures. Recent work proposed replacing constraint solvers with large language models (LLMs) to bypass these limitations, but such approaches struggle to analyze real-world codebases, where deep execution paths require globally consistent reasoning across many interacting constraints. We present \projName, a hybrid symbolic execution framework that uses LLMs selectively to generate lightweight ghost code that aids an SMT solver in handling solver-hostile code fragments, while preserving its precise, global reasoning capability. In particular, we propose three types of ghost code: (1) inversion of difficult code fragments with iterative bidirectional constraint propagation, (2) modeling via solver-friendly surrogates while preserving relevant behavior, and (3) semantic partitioning of unbounded heap spaces. We implemented \projName on top of the KLEE symbolic execution engine and evaluated it on synthetic ``logic bombs'' capturing distinct symbolic reasoning challenges, a popular mathematical library FDLibM, and \rev{four} structured-input programs (\texttt{libexpat}, \texttt{jq}, \texttt{bc} and \texttt{libyaml}). Across benchmarks, \projName{} improves coverage by \rev{28.5--115.2\%  over traditional symbolic execution baseline and by 74.1--189.8\% over LLM-based symbolic execution baselines}, while reducing LLM token usage by an average of \rev{91--96\%}. This highlights the practicality and effectiveness of this approach in real-world settings.
\end{abstract}

\maketitle


\section{Introduction}

Symbolic execution runs code on abstract inputs, encodes program behavior in logical constraints, and uses an SMT solver to analyze this behavior. The applicability of symbolic execution is hindered by several inherent limitations. First, because SMT solving is theoretically undecidable in general, symbolic execution often struggles with programs involving complex computations, as these translate into correspondingly complex SMT constraints. Second, symbolic execution requires precise modeling of program behavior, which becomes challenging in real-world environments, such as those involving network communication or extensive system calls, that are difficult to model accurately. Third, the number of possible program execution paths is effectively infinite, meaning symbolic execution can only explore a limited subset. The above challenges are particularly acute in programs that make extensive use of dynamic memory allocation and pointer-based data structures, where control flow depends on complex aliasing relations or heap configurations.

Recent developments in large language models (LLMs) promise to mitigate these limitations. In particular, Autobug~\cite{autobug}, PALM~\cite{palm} and ConcoLLMic~\cite{concollmic} use LLMs in place of, or to substantially reduce reliance on, SMT solving to directly analyze the program and generate tests for difficult program parts. While this circumvents the undecidability issues of SMT solvers, these methods struggle to analyze real-world codebases when the target functionality is deep in the program,  reachable only after several layers of calls and gating checks that require globally consistent choices about inputs and state. This is because LLMs struggle to maintain long-range, formally checkable dependencies across many functions, whereas SMT solvers accumulate and satisfy precise cross-procedural constraints, which enables reasoning about subtle interprocedural interactions.

We propose \projName, a novel way of integrating symbolic execution with LLMs. Instead of replacing SMT solvers, \projName uses an LLM only where SMT-based symbolic execution struggles. In particular, an LLM analyzes difficult parts of the program, recognizing domain structures and identifying constraint solving strategies, and translates these into lightweight ``ghost code'' that plugs into the SMT-based symbolic execution workflow as domain-specific solver aids. This preserves precise, global reasoning over cross-function dependencies and subtle interactions, while alleviating hard constraints that block progress in traditional symbolic execution.

We introduce three types of LLM-generated ghost code for helping symbolic execution analyze a complex program fragment. First, we prompt an LLM to invert the difficult fragment, so that it can be executed backward to propagate concrete values from the suffix constraints back across the fragment to satisfy prefix constraints. Second, we prompt an LLM to generate a simplified surrogate that approximates complex behavior to permit efficient forward symbolic execution. Third, for program fragments operating over dynamic data structures, we prompt an LLM to generate symbolic heap topologies based on domain-specific insights to mitigate path explosion. The choice of ghost code type is context-dependent and selected by the LLM. Soundness is preserved because any inputs or states produced with ghost code are checked against the original program, so hallucinated ghost code does not introduce unsound results.

We implemented \projName in the KLEE symbolic execution engine~\cite{klee} and compared it against previous techniques on programs that have historically challenged symbolic execution: (1) 53 synthetic programs from the LogicBombs benchmark~\cite{xu2018benchmarking} capturing distinct symbolic-reasoning challenges; (2) 78 entry points\footnote{Entry points are top-level target functions that exercise different parts of the program or library.} from the widely-used math library \texttt{FDLibM}, which induces constraints involving undecidable theories; and (3) \rev{19} source files from \rev{four} projects (\texttt{libexpat}, \texttt{jq}, \texttt{bc} and \texttt{libyaml}) targeting structured input formats and thus operating on complex dynamic data structures. Across the benchmarks, \projName{} improves coverage by \rev{28.5--115.2\% over the strongest traditional symbolic-execution baseline
and by 74.1--189.8\% over the strongest relevant LLM-based baseline} while reducing the number of tokens by \rev{91--96\%}. Via ablation, we confirmed that each type of ghost code contributed to the performance improvement.

In summary, the paper makes the following contributions:

\begin{itemize}
\item A novel enhancement of symbolic execution with LLM-generated solver-aiding ghost code.
\item Three practical ghost code variants, as well as corresponding generation prompts and solver integration algorithms, to alleviate practical challenges of symbolic execution.
\item A comprehensive evaluation on both synthetic and realistic benchmarks that highlights the strength of our approach in covering hard-to-reach program branches.
\end{itemize}

All code, data, and scripts are available at \url{https://figshare.com/s/09550df85983cbf4ae71}.

\section{Overview}

Ghost code~\cite{filliatre2016spirit} refers to auxiliary constructs added to a program to facilitate specification, reasoning, or verification without affecting its functional behavior. Traditionally, ghost code has been used in program verification to express invariants, ranking functions, or lemmas to automate proofs of program properties. Our work is the first that applies ghost code to enhance symbolic execution, and shows that it can be automatically generated using LLMs. This section illustrates the utility of LLM-generated ghost code for symbolically executing programs that involve non-trivial mathematics and dynamic data structures.

\begin{figure}
  \centering
  \begin{subfigure}[b]{0.24\textwidth}
    \centering
{
\scriptsize          
\begin{minted}[escapeinside=||]{C}
hx = __HI(x);
ix = hx & 0x7fffffff;
if (ix >= 0x7ff00000)
  return one / (x * x);
x = fabs(x);
if (ix >= 0x40000000) {
  |\hly{s = sin(x);}|
  |\hly{c = cos(x);}|
  ss = s - c;
  cc = s + c;
  if (ix < 0x7fe00000) {
    z = -cos(x + x);
    if (|\hlg{(s * c) < 0}|)
      cc = z / ss;
    else
      ss = z / cc;
    }
}
\end{minted}
}
    \caption{Function \texttt{\_ieee754\_j0} from \texttt{FDLibM} involving trigonometric operations.}
    \label{fig:ieee754}
  \end{subfigure}
  \hfill
  \begin{subfigure}[b]{0.36\textwidth}
    \centering
{
\tiny          
\begin{minted}[escapeinside=||]{C}
double golden_section(double (*f)(double, void *),
                      void *ctx,
                      double a, double b,
                      int max_iter,
                      double tol) {
  |\hlgr{algorithm omitted for brevity}|
}
                      
double sincos_error(double x, void *ctx) {
  double s = sin(x), c = cos(x);
  double ds = s - ((double *)ctx)[0];
  double dc = c - ((double *)ctx)[1];
  return ds * ds + dc * dc;
}

double f_inv(double s, double c) {
  double ctx[2] = {s, c};
  return golden_section(
    sincos_error, ctx, 0.0, 2*M_PI, 200, 1e-12
  );
}
\end{minted}
}
    \caption{LLM-generated ghost code that inverts \mintinline{C}{s = sin(x); c = cos(x);} using golden-section search algorithm.}
    \label{fig:inversion}
  \end{subfigure}
  \hfill
  \begin{subfigure}[b]{0.35\textwidth}
    \centering 
    \includesvg[width=1.4\textwidth,pretex=\relscale{0.7}]{propagation.svg}\vspace{2mm}
    \caption{sin/cos are replaced with $\alpha, \beta$ during symbolic execution; values are propagated backward by running $f^{-1}$.}
    \label{fig:propagation}
  \end{subfigure}
  \vspace{-2mm}
  \caption{\projName propagates constraints through a difficult code fragment that involves trigonometric operations using a concrete execution of an LLM-generated ghost function that inverts the fragment.}
  \label{fig:complex_math_example}
  \vspace{-2mm}
\end{figure}

\subsection{Propagating Constraints Through Difficult Code Fragments}

Symbolic execution struggles to analyze programs with complex operations such as non-linear arithmetic or trigonometric functions, since they induce hard-to-solve constraints. Consider an implementation of the Bessel function in \Cref{fig:ieee754}, which is a part of the function \texttt{\_ieee754\_j0} from the math library \texttt{FDLibM}. Suppose symbolic execution aims to explore branches guarded by the condition \mintinline[escapeinside=||]{C}`(s * c) < 0` highlighted with \hlg{\phantom{X}}. Even KLEE-Float~\cite{kleeFloat}, a version of KLEE~\cite{klee} optimized for floating-point operations, fails to solve the resulting constraints, because this check depends on the output of trigonometric operations highlighted with \hly{\phantom{X}}.

As an alternative to traditional techniques, ConcoLLMic~\cite{concollmic} replaces the SMT solver with an agentic LLM workflow to generate inputs that cover target branches. It formulates queries of the form $f(a) = b$ and then  synthesizes a suitable input $a$ using natural-language reasoning and Python tool execution. In \Cref{fig:ieee754}, although ConcoLLMic can find a satisfying assignment for constraints originating from the difficult fragment highlighted by \hly{\phantom{X}}, this fragment itself depends on the value of \mintinline[escapeinside=||]{C}`x`. The variable \mintinline[escapeinside=||]{C}`x` induces a non-trivial control-flow dependency through bit-level manipulation of its IEEE-754 representation: (1) extracting the high 32 bits with \mintinline{C}`__HI(x)`, (2) masking off the sign bit with \mintinline{C}`hx & 0x7fffffff` to obtain $|x|$, and (3) checking the guard \mintinline{C}`ix >= 0x40000000`, which corresponds to $|x| \ge 2.0$. Because ConcoLLMic lacks a mechanism for global reasoning that can combine constraints arising from multiple program locations, it fails to infer and satisfy these intertwined conditions, and therefore cannot cover the target branches guarded by \mintinline[escapeinside=||]{C}`(s * c) < 0`.

\begin{figure}
\vspace{-3mm}
  \centering
  \begin{subfigure}[b]{0.32\textwidth}
    \centering
{
\scriptsize          
\begin{minted}[escapeinside=@@]{C}
typedef struct SNode {
    int val;
    char *name;
    struct SNode *lvl0;
    struct SNode *lvl1;
} SNode;
  
int proc_if_ring(SNode *h, SNode *g) {
  SNode *a=h ? h->lvl0 : NULL;
  SNode *b=a ? a->lvl0 : NULL;
  SNode *c=b ? b->lvl0 : NULL;
  SNode *d=c ? c->lvl0 : NULL;
  if (!h||!a||!b||!c||!d) return 0;
  if (d!=h) return 0;
  if (h->lvl1!=b||a->lvl1!=c||
      b->lvl1!=d||c->lvl1!=a) return 0;
  SNode *g1 = g ? g->lvl0 : NULL;
  SNode *g2 = g1 ? g1->lvl0 : NULL;
  if (g2!=h) return 0;
  @\hlg{return proc\_ring\_values(h)}@
}
\end{minted}
}
    \caption{The function \texttt{proc\_if\_ring} checks if a skiplist forms a ring, and then processes its values.}
    \label{fig:process_ring}
  \end{subfigure}
  \hfill
  \begin{subfigure}[b]{0.28\textwidth}
    \centering
\hspace*{-2mm}\resizebox{1.2\linewidth}{!}{%
\begin{tikzpicture}[
    >=stealth,
    node distance=2cm,
    main/.style={draw,rounded corners,minimum width=11mm,minimum height=8mm,align=center},
    arr/.style={->,very thick},
    exp/.style={->,densely dashed,thick,black!70},
    gptr/.style={->,ultra thick,red!70!black,shorten >=1pt,shorten <=2pt},
    x=1cm,y=1cm
  ]
  \centering
  % --- main nodes ---
  \node[main] (n0) at (-1, 1.5) {n0};
  \node[main] (n1) at ( 1, 1.5) {n1};
  \node[main] (n2) at ( 1,-1.5) {n2};
  \node[main] (n3) at (-1,-1.5) {n3};

  % --- lvl0 ring: solid inner square ---
  \draw[arr] (n0) -- node[draw=none,fill=none,above]{lvl0} (n1);
  \draw[arr] (n1) -- node[draw=none,fill=none,right]{lvl0} (n2);
  \draw[arr] (n2) -- node[draw=none,fill=none,below]{lvl0} (n3);
  \draw[arr] (n3) -- node[draw=none,fill=none,left]{lvl0} (n0);

    \node[draw=none,fill=none,above left=15pt and 7pt of n0] (glabel) {$h$};
  \draw[gptr] (glabel.north east) -- (n0.north);

  % --- lvl1 express links: dashed, routed with outer control points ---
  \draw[exp]
    (n0) .. controls ($(n0)+(0,3.8)$) and ($(n2)+(0,3.8)$) ..
    node[draw=none,fill=none,above]{lvl1} (n2.north west);
  \draw[exp]
    (n1) .. controls ($(n1)+(2.8,0)$) and ($(n3)+(2.8,0)$) ..
    node[draw=none,fill=none,right]{lvl1} (n3.north east);
  \draw[exp]
    (n2) .. controls ($(n2)+(0,-3.8)$) and ($(n0)+(0,-3.8)$) ..
    node[draw=none,fill=none,below]{lvl1} (n0.south east);
\draw[exp]
    (n3) .. controls ($(n3)+(-2.8,0)$) and ($(n1)+(-2.8,0)$) ..
    node[draw=none,fill=none,left]{lvl1} (n1.south west);

  % --- g pointer (label only, no box) ---
  \node[draw=none,fill=none,below right=15pt and 7pt of n2] (glabel) {$g$};
  \draw[gptr] (glabel.north west) -- (n2.south);
\end{tikzpicture}
}
    \caption{Ring-shape skiplist required to reach the target function call \texttt{proc\_ring\_values}.}
    \label{fig:ring}
  \end{subfigure}
  \hfill
  \begin{subfigure}[b]{0.36\textwidth}
    \centering
{
\tiny          
\begin{minted}[escapeinside=!!]{C}
void constrain_heap(SNode *n0, SNode *n1,
                    SNode *n2, SNode *n3,
                    SNode **h, SNode **g) {
  switch (!$\xi$!) {
    case 1: /* broken chain */
      n0->lvl0 = n1; n1->lvl0 = NULL;
      n0->lvl1 = NULL; n1->lvl1 = NULL;
      n2->lvl1 = NULL; n3->lvl1 = NULL;
      *h = n0; *g = n3;
      break;

    !\hlgr{three cases omitted for brevity}!

    case 5: /* ring + g two-steps to h */
      n0->lvl0 = n1; n1->lvl0 = n2;
      n2->lvl0 = n3; n3->lvl0 = n0;
      n0->lvl1 = n2; n1->lvl1 = n3;
      n2->lvl1 = n0; n3->lvl1 = n1;
      *h = n0;
      *g = n2;
      break;
  }
}
\end{minted}
}
    \caption{LLM-generated ghost code that constrains heap with five semantic topologies relevant to the problem context.}
    \label{fig:constrain_heap}
  \end{subfigure}
  \vspace{-2mm}
  \caption{\projName tackles heap-configuration explosion by using LLM-generated ghost code to constrain the heap to semantically meaningful data-structure shapes while preserving symbolic contents.}
  \label{fig:topology_example}
  \vspace{-2mm}  
\end{figure}

To address this challenge, \projName operates in three steps. In step 1, it asks an LLM to identify a code fragment that is hard for an SMT solver to reason about and to generate ghost code that inverts the fragment. In this example, the LLM flags \mintinline{C}`s = sin(x); c = cos(x);`, which we refer to as $f$, as difficult and produces an efficient inverse procedure \(f^{-1}\) using golden-section search~\cite{kiefer1953sequential} as shown in \Cref{fig:inversion}. In step 2, \projName executes the program on a symbolic input \(x\), replacing the difficult fragment with fresh symbolic variables \(\alpha\) and \(\beta\) that stand for the values of \(s\) and \(c\) as shown in \Cref{fig:propagation}. During this run it collects two constraints: (1) a prefix path condition \(\pi_{\mathrm{pre}}\) from code preceding the fragment, which depends only on the input \(x\); and (2) a suffix path condition \(\pi_{\mathrm{suf}}\) from code following the fragment, which depends on \(\alpha\) and \(\beta\).

The final step integrates the synthesized ghost code into the constraint-solving workflow via \emph{bidirectional constraint propagation}. First, \projName solves the suffix constraints \(\pi_{\mathrm{suf}}\) with an SMT solver to obtain candidate values for \(s\) and \(c\), denoted \(\alpha_0\) and \(\beta_0\); assume the solver returns \((\alpha_0,\beta_0)=(10,-1)\). Next, \projName propagates these values backward through the difficult fragment by concretely executing the synthesized inverse \(f^{-1}\). Although no real \(x\) can satisfy these values simultaneously (since \(\sin^2(x)+\cos^2(x)=1\)), the inverse computation yields the closest feasible input \(x_0 \approx 1.67\) radians. This input violates the prefix constraint \(\pi_{\mathrm{pre}}\), which requires \(|x|\ge 2.0\). \projName therefore uses an optimization solver to find the closest \(x_1\) to \(x_0\) that satisfies \(\pi_{\mathrm{pre}}\); in this example, \(x_1=2.0\). It then propagates \(x_1\) forward by executing the original fragment \(f\), producing \((\alpha_1,\beta_1)\approx(0.909,-0.416)\), which satisfies \(\pi_{\mathrm{suf}}\). Finally, \projName concretely executes the program to validate that the generated input indeed reaches the target branch, mitigating potential ghost code inaccuracies due to LLM hallucinations.

In comparison to traditional symbolic execution, \projName takes advantage of LLM's domain knowledge and code generation capabilities in the form of an inversion algorithm, for which it chose to implement a golden-section search. In comparison to LLM-based techniques such as ConcoLLMic, \projName takes advantage of the SMT solver's precise reasoning about cross-function interactions, i.e. \mintinline{C}{_ieee754_j0} and \mintinline{C}{__HI}, achieved through our bidirectional constraint propagation algorithm that plugs the generated ghost code into the solver logic. Through this synergy, \projName is the only tool that provides full coverage for this function.

\subsection{Constraining Heap With Semantic Topologies}

A major limitation of symbolic execution in real-world testing is heap-configuration explosion caused by dynamic data structures, such as structured-input parse trees and the complex list/tree structures used in databases and file systems. Due to the complexity of these data structures and their processing code, we demonstrate how \projName addresses this problem on a synthetic example of a simplified 2-level skiplist shown in \Cref{fig:process_ring}. In this data structure, \mintinline{C}`lvl0` points to the next node in the full (level-0) linked list, while \mintinline{C}`lvl1` is an ``express lane'' pointer that skips over some nodes to speed up searches. Each node stores a key/value \mintinline{C}`val` and associated data \mintinline{C}`name`.

Let the function under test be \mintinline{C}{proc_if_ring} in \Cref{fig:process_ring}. It checks whether a given skiplist forms a valid ring, as in \Cref{fig:ring}, and, if so, processes the values stored in its nodes. Testing this function is challenging because reaching the call to the value-processing routine (highlighted with \hlg{\phantom{X}}) requires symbolic execution to explore a prohibitive number of heap configurations, covering many possible shapes and aliasing relationships among ten pointers. As a result, KLEE fails to reach this function within the timeout of 10 minutes. On the other hand, the LLM-based tool Cottontail~\cite{cottontail} accounts for input structure and can generate valid data-structure shapes, but --- similarly to ConcoLLMic --- it cannot reconcile value constraints originating elsewhere (e.g., in \texttt{proc\_ring\_values}), limiting its effectiveness on large projects.

To address these limitations, \projName leverages an LLM's contextual understanding to generate ghost code, \texttt{constrain\_heap}, which constrains the heap with \emph{semantic topologies}. A semantic topology is a partially concrete, partially symbolic heap configuration that captures a meaningful use case drawn from the problem domain or program context. In this example, the LLM generates five semantic topologies --- chain, non-ring, misaligned express pointers, mismatched head link, and a valid ring -- shown in \Cref{fig:constrain_heap}. The selected topology is determined by a fresh symbolic variable $\xi$. When symbolically executing the program, \projName inserts a call of \texttt{constrain\_heap} in front of each call of \mintinline{C}{proc_if_ring}:

{\scriptsize
\begin{minted}[escapeinside=@@]{C}
/* n0, n1, n2, n3, h, g are initialized with symbolic memory */
constrain_heap(n0, n1, n2, n3, h, g);
int result = proc_if_ring(h, g);
\end{minted}
}

\noindent Crucially, within each topology, only the links relevant to the intended shape are fixed; all other fields remain symbolic, including the value and name fields. This enables exploration of shape-dependent branches without over-constraining subsequent execution. As a result, \projName efficiently covers all branches in the function under test, as well as the function \texttt{proc\_ring\_values}, within 2 seconds.

\section{Background \& Notation}

We work in the standard Satisfiability Modulo Theories (SMT) setting. Let $X$ denote a finite set of (typed) variables, ranged over by $\alpha,\beta,\gamma,\xi\in X$, and let there be a fixed signature (function and predicate symbols) interpreted in a background theory $T$ (e.g., linear arithmetic).
We write $\Phi_{T}(X)$, or just $\Phi$ for short, for the set of all well-formed formulas with free variables in $X$; formulas are denoted by $\phi,\psi\in \Phi$.
An assignment (or model candidate) is a map $A: X \to \mathrm{Dom}(T)$; satisfaction is written $A \models_T \phi$.
A formula $\phi$ is $T$-satisfiable if there exists $A$ such that $A \models_T \phi$, and $T$-unsatisfiable otherwise.
We denote by $\mathrm{SMT}(\phi)$ an SMT-solver call that returns either \textsc{sat} together with a satisfying assignment/model $A$, or \textsc{unsat}.
For optimization, let $f$ be a (theory) term denoting an objective value (typically numeric). We write $\mathrm{OPT}(f,\phi)$ for an optimization-modulo-theories (OMT) call returning an assignment $A^\star$ such that $A^\star \models_T \phi$ and $f(A^\star)$ is optimal among all $A$ satisfying $\phi$ (with \textsc{min} or \textsc{max} understood from context), or returning \textsc{unsat} if $\phi$ is unsatisfiable.
For substitution, given a term $t$ and variable $x\in X$, we write $\phi[x\mapsto t]$ for the capture-avoiding substitution of $t$ for all free occurrences of $x$ in $\phi$.

Let $\mathcal{P}$ denote the set of programs, ranged over by $P,Q\in\mathcal{P}$, and let $\mathcal{S}$ denote the set of statements, ranged over by $s,t\in\mathcal{S}$; programs are (finite) sequences of statements, written $P = s_1;\cdots;s_n$. We assume a set of program variables $\mathsf{Var}$, denoted by $v,w$, distinct from the set of logical/symbolic variables $X$ used in formulas/constraints. Procedure calls have the form $\mathsf{call}\;p(e_1, ..., e_n)$ for $p\in\mathsf{Proc}$ and argument expressions $e_1, ..., e_n$; we write $P[C \mapsto \mathsf{call}\;p(e_1, ..., e_n)]$ for replacing a designated code fragment (subprogram) $C$ occurring in $P$ by that call.

A (concrete) program state is a pair $\sigma=(\nu,h)$ where the store $\nu:\mathsf{Var}\to\mathrm{Val}$ maps program variables to concrete values and the heap $h:\mathrm{Addr}\rightarrow \mathrm{Cell}$ is a finite partial map from addresses to heap cells (records), with $\mathrm{Cell}$ determined by the language (e.g., tuples of fields); a symbolic state is a triple $\hat\sigma=(\hat\nu,\hat h,\pi)$ where the symbolic store $\hat\nu:\mathsf{Var}\to\mathrm{Term}(X)$ maps variables to SMT terms, the symbolic heap $\hat h:\mathrm{Addr}\rightarrow \mathrm{CellTerm}(X)$ maps addresses to (possibly symbolic) cells whose fields are SMT terms, and $\pi\in\Phi_{T}(X)$ is the path condition constraining the symbols. All concrete states are denoted as $\Sigma$, and symbolic states as $\widehat{\Sigma}$. Satisfaction/realization is given by an assignment $A$ such that $A\models_T \pi$, yielding a concrete state $\sigma_A=(\nu_A,h_A)$ by evaluating all terms in $\hat\nu$ and $\hat h$ under $A$ (writing $\llbracket t\rrbracket_A$ for term evaluation). A concrete program semantics is a transition relation $\to \;\subseteq\; (\mathcal S\times \Sigma)\times(\mathcal S\times \Sigma)$, and its reflexive–transitive closure is written $\to^{*}$. Then, a concrete execution function is defined as $\mathrm{exec}(P,\sigma)\;=\;\sigma' \in \Sigma \text{ such that } \langle P,\sigma\rangle \to^{*} \langle \_,\sigma'\rangle.$

\begin{example}[Concrete execution of a simple assignment]
Let $P$ be the program $x := x + 1$ and let the initial store be
$\sigma$ with $\sigma(x)=41$. Since $\llbracket x+1\rrbracket_\sigma = 42$, we have $\langle x := x + 1,\sigma\rangle \to \langle \_, \sigma[x \mapsto 42]\rangle$. Therefore, $\mathrm{exec}(x := x+1,\sigma)=\sigma[x \mapsto 42]$.
\end{example}

A symbolic semantics is a transition relation $\Rightarrow \;\subseteq\; (\mathcal S\times \widehat\Sigma)\times(\mathcal S\times \widehat\Sigma)$ that updates symbolic store/heap with symbolic expressions and updates (conjoins) the path condition, while discarding unsat branches. Then, we define symbolic execution as a function $\mathrm{symex}:\mathcal P\times \widehat\Sigma \to 2^{\widehat\Sigma}$ such that
\[
\mathrm{symex}(P,\hat\sigma)\;=\;\{\hat\sigma' \in \widehat\Sigma \mid \langle P,\hat\sigma\rangle \Rightarrow^{*} \langle \_,\hat\sigma'\rangle\}.
\]

\begin{example}[Symbolic execution with a branch]
Let the program be $P \triangleq \texttt{if }(x > 0)\texttt{ then }y := x+1\texttt{ else }y := 0$. Assume we start from the symbolic state
\[
\hat\sigma_0=(\hat\nu_0,\hat h_0,\pi_0)
\quad\text{with}\quad
\hat\nu_0(x)=\alpha,\;\hat\nu_0(y)=\beta,\;\hat h_0=\emptyset,\;\pi_0=\top,
\]
where $\alpha,\beta\in X$ are fresh logical variables. Symbolic execution explores both branches, adding the corresponding guard to the path condition $\langle P,\hat\sigma_0\rangle \Rightarrow^{*} \langle \_,\hat\sigma_1\rangle$ and $\langle P,\hat\sigma_0\rangle \Rightarrow^{*} \langle \_,\hat\sigma_2\rangle$, where
\begin{align*}
&\hat\sigma_1 = (\hat\nu_1,\emptyset,\pi_1),
\quad
\pi_1 \equiv \alpha>0,
\quad
\hat\nu_1(y)=\alpha+1,\;\hat\nu_1(x)=\alpha,\\
&\hat\sigma_2 = (\hat\nu_2,\emptyset,\pi_2),
\quad
\pi_2 \equiv \alpha\le 0,
\quad
\hat\nu_2(y)=0,\;\hat\nu_2(x)=\alpha.
\end{align*}
Thus, $\mathrm{symex}(P,\hat\sigma_0)=\{\hat\sigma_1,\hat\sigma_2\}$.
\end{example}

\section{\projName}

\begin{figure}[t]
    \centering
        {\scriptsize
      \begin{tcolorbox}[mybox,title={Ghost Inversion Prompt Template}]
        \begin{minted}[escapeinside=||]{text}
Generate a function {SIGN} that inverts the original code fragment to recover the pre-state values of the
variables {VARS}. If exact inversion is not possible, implement a bounded numeric search or a clamped
approximation, and clearly comment on this in the output. The original function code: {CODE}
      \end{minted}
      \end{tcolorbox}
        }
\vspace{-2mm}
\caption{\label{fig:inverse_template}A simplified fragment of \projName's prompt template for generating ghost code that inverts a difficult fragment, where \texttt{SIGN} is the signature of the desired function, \texttt{VARS} are the variables, whose values the inversion needs to compute, and \texttt{CODE} is the code fragment.}
\vspace{-2mm}
\end{figure}
\projName{} adds a layer on top of a traditional SMT-backed symbolic execution engine. 
\rev{\projName{} first performs an exploratory KLEE run with solver timeouts reduced by 90\% to identify uncovered branches, which mark hard-to-reach paths. These candidate branches are then ranked by how high they occur in the control flow and by the LOC size of the region they guard, prioritizing branches whose resolution could unlock larger still-uncovered code regions. For each branch, \projName{} constructs the LLM input \{CODE\} from the branch condition, the enclosing function, preceding code that may influence it, and directly related callees, truncated to a fixed context budget.} Based on this fragment
and the LLM's choice, \projName{} generates one of three kinds of ghost code to improve constraint solving and path exploration: (1) fragment inversion, (2) surrogate models, and (3) semantic heap partitioning. 
\rev{The process is fully automated, not requiring manual hints or hand-crafted inversions.} 
The ghost code is then integrated into the symbolic executor and solver as explained below.

\subsection{Bidirectional Constraint Propagation with Code Inversion}
\label{sec:bidirectional}

\begin{figure}[t]
  \begin{subfigure}[b]{0.48\textwidth}
    \centering
{\scriptsize          
\begin{minted}[escapeinside=||]{C}
int hard_func(int x) { 
    return (x * x + 6) % 97; 
}
int main(void) {
    int a = |$\alpha$|;
    |\hly{int b = hard\_func(a);}|
    if (b == 42) {
        assert(a > 20);
    }
}
\end{minted}
}
  \end{subfigure}
  \hfill
  \begin{subfigure}[b]{0.48\textwidth}
    \centering
{\scriptsize          
\begin{minted}[escapeinside=||]{C}
void f_inv(int b, int *results) {
    int count = 0;
    int t = (b - 6) % 97;
    if (t < 0) t += 97;
    for (int a = 0; a < 97; a++) {
        if ((a * a) % 97 == t) {
            results[count++] = a;
            if (count == 2) break;
        }
    }
}
\end{minted}
}
  \end{subfigure}
  \vspace{-2mm}
  \caption{Running example to illustrate ghost inversion and bidirectional constraint propagation.\label{fig:inversion_running_example}} 
  \vspace{-2mm} 
\end{figure}


Given a program $P$ in which an LLM has identified a fragment $C$ as difficult for SMT reasoning, e.g. non-linear arithmetic, \projName uses an LLM to generate a ghost inversion procedure. Intuitively, we treat $C$ as a (partial) state transformer on a selected slice of the state --- i.e., the set of program variables and heap locations that $C$ may read or write --- and synthesize code that implements an inverse of that transformer. Executing the ghost inversion in reverse then propagates constraints from the post-state back to the pre-state, allowing the solver to reason about $C$ indirectly even when its forward semantics is not handled well by the underlying theory.

Let the sets of program variables and heap locations that a code fragment \(C\) may read/write be
\[
\mathsf{RdVar}(C),\mathsf{WrVar}(C)\subseteq \mathsf{Var},\qquad
\mathsf{RdLoc}(C),\mathsf{WrLoc}(C)\subseteq \mathrm{Addr}.
\]
Then, the read/write footprints of \(C\) are defined as
\[
\mathsf{Rd}(C)\;\triangleq\;\mathsf{RdVar}(C)\cup \mathsf{RdLoc}(C),
\qquad
\mathsf{Wr}(C)\;\triangleq\;\mathsf{WrVar}(C)\cup \mathsf{WrLoc}(C),
\]
and we denote \(\sigma\!\upharpoonright_{S}\) for projecting a concrete state \(\sigma\) to the values of a set \(S\) of variables/heap locations. We model \(C\) as a function from the values $C$ reads to the values it writes, $f_C \;:\; \Sigma_{\mathsf{Rd}(C)} \rightarrow \Sigma_{\mathsf{Wr}(C)}$:
\[
f_C(\rho)=\omega
\;\;\triangleq\;\;
\exists \sigma.\;
\Bigl(\rho=\sigma\!\upharpoonright_{\mathsf{Rd}(C)}
\;\wedge\;
\omega=\mathrm{exec}(C,\sigma)\!\upharpoonright_{\mathsf{Wr}(C)}\Bigr),
\]
\noindent Using the prompt shown in \Cref{fig:inverse_template}, the LLM is requested to generate a ghost function \mintinline{C}`f\_inv` that implements an inverse $f^{-1}_C \;:\; \Sigma_{\mathsf{Wr}(C)} \rightharpoonup \Sigma_{\mathsf{Rd}(C)}$ such that
\[
\forall \rho\in \mathrm{Dom}(f_C).\;\;
f^{-1}_C\!\bigl(f_C(\rho)\bigr)=\rho.
\]
Operationally, this means that for any concrete state \(\sigma\) where both executions terminate,
\[
\mathrm{exec}(\mathtt{f\_inv},\,\mathrm{exec}(C,\sigma))\!\upharpoonright_{\mathsf{Rd}(C)}
\;=\;
\sigma\!\upharpoonright_{\mathsf{Rd}(C)}.
\]

In practice, many fragments cannot be fully inverted, and an LLM is instead requested to return a set of possible values that are all valid inverses, i.e. $\forall \rho\in \mathrm{Dom}(f_C).\;\; \rho \in f^{-1}_C\!\bigl(f_C(\rho)\bigr)$, or when no inverse values exist, the generated function is tasked to find an input value that results in the closest output state to the expected state, as shown in the example in \Cref{fig:inversion}.

\begin{example}[Ghost inversion for a hard fragment]
  Consider the program fragment in \Cref{fig:inversion_running_example} (left). Let $C$ be the call/statement $b := \mathsf{hard\_func}(a)$. Its read/write sets are 
\[
\mathsf{RdVar}(C)=\{a\},\qquad \mathsf{WrVar}(C)=\{b\},\qquad \mathsf{RdLoc}(C)=\mathsf{WrLoc}(C)=\emptyset,
\]
hence $\mathsf{Rd}(C)=\{a\}$ and $\mathsf{Wr}(C)=\{b\}$. For this fragment, an LLM generates the function \mintinline{C}`f_inv` shown on the right of the same figure, where the output of the function is the list of values \mintinline{C}`results`.
\end{example}
  
\projName replaces $C$ with an assignment of $\mathsf{Wr}(C)$ to fresh symbolic variables, since it is not able to directly reason about the output of $C$ using an SMT solver. Assume \(\mathsf{Wr}(C)=\{\nu_1,\dots,\nu_k\}\), and \(\beta_1,\dots,\beta_k\in X\) be fresh logical symbols w.r.t. the current symbolic state. Let the havoc statement be
\[
\mathsf{havoc}(\mathsf{Wr}(C)) \;\triangleq\; (\nu_1 := \beta_1;\;\cdots;\; \nu_k := \beta_k).
\]
Then, \projName transforms the program by replacing $C$ with this havoc statement:
\[
P' \triangleq P[C \mapsto \mathsf{havoc}(\mathsf{Wr}(C))].
\]

Finally, \projName symbolically executes $P'$ from an initial symbolic state
$\hat\sigma_0=(\hat\nu_0,\hat h_0,\pi_0)$, which is $\hat h_0=\emptyset$, $\pi_0\equiv \top$, and $\hat\nu_0(a)=\alpha$ (fresh $\alpha\in X$) in the running example. Let the set of final symbolic states produced by symbolic execution of the transformed program be
\[
\mathrm{symex}(P',\hat\sigma_0)
\;=\;
\{\hat\sigma_1^1,\ldots,\hat\sigma_1^m\},
\]

\noindent where each \(\hat\sigma_1^i=(\hat\nu_1^i,\hat h_1^i,\pi_1^i)\) is a (feasible) terminal symbolic state. To make explicit what happens before and after the replaced fragment \(C\), we decompose each $\pi_1^i$ into corresponding prefix and suffix constraints: $\pi_1^i = \pi_\mathrm{pre}^i \wedge \pi_\mathrm{suf}^i$, where $\pi_\mathrm{pre}^i$ depends only on the initial input states, while $\pi_\mathrm{suf}^i$ also depends on the additional symbolic variables introduced by the havoc statement\footnote{The path condition has to be decomposed into a sequence of sub-constraints if $C$ is executed inside a loop.}. In the running example, the two final path conditions decompose as:
\begin{align*}
\pi_f^{(1)} \equiv (\,\beta = 42\,)\wedge(\,\alpha>20\,)
\;&=\;
\underbrace{\top}_{\pi_{\mathrm{pre}}^{(1)}}\wedge
\underbrace{((\beta=42)\wedge(\alpha>20))}_{\pi_{\mathrm{suf}}^{(1)}},\\
\pi_f^{(2)} \equiv (\,\beta \neq 42\,)
\;&=\;
\underbrace{\top}_{\pi_{\mathrm{pre}}^{(2)}}\wedge
\underbrace{(\beta\neq 42)}_{\pi_{\mathrm{suf}}^{(2)}}.
\end{align*}

For each pair of $\pi_\mathrm{pre}^i, \pi_\mathrm{suf}^i$, after computing the above path conditions, \projName executes our bidirectional constraint propagation algorithm (\Cref{alg:bidirectional_reconcile}) that reconciles prefix and suffix constraints around a transformed fragment \(f\) by iteratively ``meeting in the middle'' using the generated inverse \(f^{-1}\). It first solves the current suffix constraint \(\pi_{\mathrm{suf}}\) to obtain an initial assignment \(A_{\mathrm{suf}}\); from this it extracts proposed values for the variables written by \(f\) and runs \(f^{-1}\) to propagate them backward to a proposed pre-state \(\sigma_{\mathrm{pre}}\) over the variables read by \(f\). It then finds an assignment \(A_{\mathrm{pre}}\) satisfying \(\pi_{\mathrm{pre}}\) that is as close as possible (under a domain-specific distance discussed below) to \(\sigma_{\mathrm{pre}}\), propagates \(A_{\mathrm{pre}}\) forward through \(f\) to get an output state \(\sigma_{\mathrm{out}}\), and symmetrically optimizes within \(\pi_{\mathrm{suf}}\) to update \(A_{\mathrm{suf}}\) so that it is close to \(\sigma_{\mathrm{out}}\). If it satisfies \(\pi_{\mathrm{pre}}\wedge\pi_{\mathrm{suf}}\), it returns it; otherwise it repeats the loop until it reaches a fixpoint or a solver call on \(\pi_{\mathrm{pre}}\) or \(\pi_{\mathrm{suf}}\) returns \textsc{Unsat}. 
\rev{Under an approximate inverse, \textsc{Unsat} is operational: \projName{} found no satisfying assignment within the current approximation and budget; the path itself may or may not be feasible.}

The algorithm drives the prefix--suffix reconciliation toward a fixpoint by treating the values propagated across the cut as soft targets and enforcing consistency via optimization. Concretely, after executing \(f^{-1}\) (resp., \(f\)) we obtain a suggested valuation \(\sigma_{\mathrm{pre}}\) for the pre-variables (resp., \(\sigma_{\mathrm{out}}\) for the post-variables). We then choose, among all assignments satisfying the hard constraints \(\pi_{\mathrm{pre}}\) (resp., \(\pi_{\mathrm{suf}}\)), one that stays as close as possible to the suggested valuation.


When the interface between the two sides is a single numeric variable \(x\), we use Z3’s optimization engine~\cite{bjorner2015nuz} and minimize an \(\ell_1\) distance objective:
\[
A_{\mathrm{pre}} \;=\; \arg\min_{A \models \pi_{\mathrm{pre}}}\; |A(x)-\sigma_{\mathrm{pre}}(x)|,
\qquad
A_{\mathrm{suf}} \;=\; \arg\min_{A \models \pi_{\mathrm{suf}}}\; |A(x)-\sigma_{\mathrm{out}}(x)|.
\]
This makes each update the smallest possible correction needed to restore feasibility, which empirically stabilizes the iteration. When the interface consists of multiple variables \(V=\{x_1,\ldots,x_n\}\), optimizing a numeric distance can scale poorly. Instead, we bias updates toward changing as few variables as possible. We encode, for each \(x_i \in V\), a soft equality \(x_i=\hat x_i\) where \(\hat x_i\) is the suggested value from \(\sigma_{\mathrm{pre}}\) (or \(\sigma_{\mathrm{out}}\)), and solve a \emph{partial MaxSAT} instance with \(\pi_{\mathrm{pre}}\) (or \(\pi_{\mathrm{suf}}\)) as hard constraints:
\[
\max\;\sum_{i=1}^n \mathbf{1}\bigl(A(x_i)=\hat x_i\bigr)
\quad\text{s.t.}\quad A \models \pi_{\mathrm{pre}} \;\;(\text{resp. } \pi_{\mathrm{suf}}).
\]

\begin{wrapfigure}{r}{0.5\textwidth}
\resizebox{0.9\textwidth}{!}{%
  \begin{minipage}{\textwidth}
  \begin{algorithm}[H]
\KwIn{Prefix constraint $\pi_{\mathrm{pre}}$, suffix constraint $\pi_{\mathrm{suf}}$,\\difficult fragment $f$, synthesized inverse $f^{-1}$}
\KwOut{A model/assignment $A$ or \textsc{Unsat}}

\BlankLine
$\pi \gets \pi_{\mathrm{pre}} \wedge \pi_{\mathrm{suf}}$\;
$\mathsf{target}_{\mathrm{suf}} \gets \mathrm{vars}(\pi_{\mathrm{suf}})$\;
$\mathsf{target}_{\mathrm{pre}} \gets \mathrm{vars}(\pi_{\mathrm{pre}})$\;
$A_{\mathrm{suf}} \gets \mathrm{SMT}\bigl(\pi_{\mathrm{suf}}\bigr)$\;
\If{$A_{\mathrm{suf}}=\bot$}{\Return{\textsc{Unsat}}}

\BlankLine
\While{\textbf{true}}{

  $\sigma_{\mathrm{pre}} \gets \mathrm{exec}\bigl(f^{-1},\, A_{\mathrm{suf}}\!\upharpoonright_{\mathsf{Wr}(f)}\bigr)$\;

  $A_{\mathrm{pre}} \gets \mathrm{OPT}\Bigl(
      \mathrm{dist}\bigl(\mathsf{target}_{\mathrm{pre}},\, \sigma_{\mathrm{pre}}\bigr),
      \ \pi_{\mathrm{pre}}
    \Bigr)$\;
  \If{$A_{\mathrm{pre}}=\bot$}{\Return{\textsc{Unsat}}}

  $\sigma_{\mathrm{out}} \gets \mathrm{exec}\bigl(f,\, A_{\mathrm{pre}}\!\upharpoonright_{\mathsf{Rd}(f)}\bigr)$\;

  $A_{\mathrm{suf}} \gets \mathrm{OPT}\Bigl(
      \mathrm{dist}\bigl(\mathsf{target}_{\mathrm{suf}},\, \sigma_{\mathrm{out}}\bigr),
      \ \pi_{\mathrm{suf}}
    \Bigr)$\;
  \If{$A_{\mathrm{suf}}=\bot$}{\Return{\textsc{Unsat}}}

  \If{$A_{\mathrm{suf}} \models \pi_{\mathrm{pre}} \wedge \pi_{\mathrm{suf}}$}{
    \Return{$A_{\mathrm{suf}}$}
  }
}
  \end{algorithm}
  \end{minipage}
}
\vspace{-3mm}
\caption{Bidirectional constraint propagation.\label{alg:bidirectional_reconcile}}
\vspace{-3mm}
\end{wrapfigure}


In our running example, \projName uses the ghost inverse \mintinline{C}{f_inv} to enumerate candidate pre-values consistent with the observed havoc value $\{\,a \in \mathrm{Val} \mid a \in f^{-1}_C(42)\,\}=\{6,91\}.$ Finally, candidates are filtered by the remaining suffix (post-fragment) constraints over pre-variables (here, $\alpha>20$), leaving $a=91$ as the only value that satisfies the explored path. \Cref{fig:complex_math_example} shows a more complex example.

\begin{figure}[t]
    \centering
        {\scriptsize
      \begin{tcolorbox}[mybox,title={Ghost Model Prompt Template}]
        \begin{minted}[escapeinside=||]{text}
Rewrite the original fragment so that (1) functionality (outputs, branches) is preserved, (2) the rewritten
code translates to a solver-friendly SMT theory such as bitvectors or linear arithmetic. For any external
library or crypto call, provide a stub implementation. The original code fragment is {CODE}.
      \end{minted}
      \end{tcolorbox}
        }
\vspace{-3mm}
\caption{\label{fig:model_template}A fragment of \projName's prompt template for generating ghost model for a difficult fragment in \texttt{CODE}.}
\vspace{-4mm}
\end{figure}

\subsection{Surrogate Models}
\label{sec:surrogate-models}

Some difficult fragments are poor candidates for inversion: they may discard information (many-to-one), or depend on external state. In such cases, \projName asks an LLM to synthesize a \emph{surrogate model} --- a solver-friendly replacement that preserves the fragment's relevant behavior while avoiding solver-hostile theories.
Let \(C\) be the difficult fragment, with read/write footprints
\(\mathsf{Rd}(C)\) and \(\mathsf{Wr}(C)\), modeled as a state transformer
\(f_C:\Sigma_{\mathsf{Rd}(C)} \to \Sigma_{\mathsf{Wr}(C)}\).
Using the prompt template in \Cref{fig:model_template}, \projName{} queries an
LLM to generate a ghost fragment \(C_{\mathsf{sur}}\) implementing a surrogate
transformer \(g_C:\Sigma_{\mathsf{Rd}(C)} \to \Sigma_{\mathsf{Wr}(C)}\), which
approximates \(f_C\) while yielding solver-friendly constraints.
The program is rewritten as:
\[
P' \triangleq P[C \mapsto C_{\mathsf{sur}}],
\qquad
\mathrm{symex}(P',\hat\sigma_0)=\{\hat\sigma_1^1,\ldots,\hat\sigma_1^m\}.
\]
Each terminal symbolic state
\(\hat\sigma_1^i=(\hat\nu_1^i,\hat h_1^i,\pi_1^i)\) induces a simplified path
condition \(\pi_1^i\) corresponding to the original program \(P\).
Because \(g_C\) may be approximate, \projName treats it as a heuristic aid rather than a trusted semantic replacement. Any concrete input produced by solving constraints from \(P'\) is validated by re-executing the original program \(P\) to ensure the intended path is feasible in \(P\). \rev{This guarantees soundness of reported coverage but not completeness: approximate inversions, surrogate imprecision or early non-convergence may cause \projName{} to miss feasible paths.}

\begin{figure}[t]
    \centering
        {\scriptsize
      \begin{tcolorbox}[mybox,title={Ghost Topology Builder Prompt Template}]
        \begin{minted}[escapeinside=||]{text}
Identify a small number of distinct heap or pointer configurations that drive different control-flow outcomes.
Generate a function {SIGN} that builds representative memory shapes and pointer configurations. Leave concrete
values in the data structures symbolic. Code context: {CODE}
      \end{minted}
      \end{tcolorbox}
        }
\vspace{-3mm}
\caption{\label{fig:topology_template}A fragment of \projName's prompt template for generating ghost code that constrains heap, where \texttt{SIGN} is the expected function signature, \texttt{CODE} is a heap manipulating function, and data structure definitions.}
\vspace{-3mm}
\end{figure}

\subsection{Semantic Heap Partitioning}

Symbolic execution often fails on pointer-rich code because many possible aliasing patterns and heap shapes cause path explosion. \projName mitigates this by generating ghost code that partitions the unbounded heap space into a small set of \emph{semantic topologies} --- partially concrete, partially symbolic heap configurations corresponding to meaningful shapes in the program's domain. Such ghost code is generated automatically for all functions that operate on unbounded symbolic pointers.

\begin{example}[List segment of length $\ge 3$.]\label{example:topology}
Let $\mathsf{head},\mathsf{n}_1,\mathsf{n}_2 \in \mathsf{Var}$ be program variables of pointer type, and $\mathsf{next}$ and $\mathsf{val}$ be data fields. Let $X$ contain address-typed symbols $\alpha_0,\alpha_1,\alpha_2,\alpha_3$ and value-typed symbols $\beta_0,\beta_1,\beta_2$. We define a \emph{semantic topology} that represents the shape
\[
\alpha_0 \xrightarrow{\mathsf{next}} \alpha_1 \xrightarrow{\mathsf{next}} \alpha_2 \xrightarrow{\mathsf{next}} \alpha_3
\]
where $\alpha_0,\alpha_1,\alpha_2$ are the first three concrete links, node values $\beta_0,\beta_1,\beta_2$ are symbolic, and the tail pointer $\alpha_3$ remains symbolic (so the list has length $\ge 3$, with an unconstrained remainder). This semantic topology is represented via the symbolic state $\hat\sigma=(\hat\nu,\hat h,\pi)$, where the first three nodes are \emph{concretely} linked, while the tail pointer is symbolic:
\begin{align*}
&\hat\nu(\mathsf{head}) = \alpha_0,\qquad
\hat\nu(\mathsf{n}_1) = \alpha_1,\qquad
\hat\nu(\mathsf{n}_2) = \alpha_2, \qquad \hat h(\alpha_0) = \langle \mathsf{val} \mapsto \beta_0,\;\mathsf{next} \mapsto \alpha_1\rangle,\\
&\hat h(\alpha_1) = \langle \mathsf{val} \mapsto \beta_1,\;\mathsf{next} \mapsto \alpha_2\rangle,\quad
\hat h(\alpha_2) = \langle \mathsf{val} \mapsto \beta_2,\;\mathsf{next} \mapsto \alpha_3\rangle.
\end{align*}
The path condition $\pi$ enforces that the three nodes are distinct, allocated, and that $\alpha_3$ denotes an arbitrary (possibly null) tail:
\[
\pi \triangleq (\alpha_0 \neq \mathsf{null}) \wedge (\alpha_1 \neq \mathsf{null}) \wedge (\alpha_2 \neq \mathsf{null})
\wedge (\alpha_0 \neq \alpha_1) \wedge (\alpha_1 \neq \alpha_2) \wedge (\alpha_0 \neq \alpha_2).
\]
\end{example}

\projName adds heap-topology constraints by executing generated ghost code. The key idea is to (i) introduce
a fresh selector variable $\xi$ and (ii) conditionally constrain the symbolic
heap to satisfy one of a small set of semantic topologies identified based on the definition of the data structure, and the code context. Each topology yields a solver-friendly, bounded set of constraints.

% Let the fragment $C$ be a call site $\mathsf{call}\;p(r)$ where the argument $r$ is pointer-typed and fully symbolic. Using the prompt template in \Cref{fig:topology_template}, \projName asks the LLM to synthesize a ghost procedure $\mathsf{constrain\_heap}_p(r,\xi)$, where $\xi\in X$ is a fresh symbolic variable that selects among candidate topologies. The original program is rewritten by inserting the ghost call immediately before the call $C$:
% \[
% P' \triangleq P[C \mapsto \mathsf{constrain\_heap}_p(r,\xi);\;C].
% \]
Let the fragment \(C\) be a call site \(\mathsf{call}\;p(r)\) where the argument
\(r\) is pointer-typed and fully symbolic. Using the prompt template in
\Cref{fig:topology_template}, \projName{} asks the LLM to synthesize a ghost
procedure \(\mathsf{constrain\_heap}_p(r,\xi)\), where \(\xi \in X\) is a fresh
symbolic variable selecting among candidate heap topologies. The program is then
rewritten as \(P' \triangleq P[C \mapsto \mathsf{constrain\_heap}_p(r,\xi);\,C]\),
inserting the ghost call immediately before \(C\).
Symbolic execution then proceeds on $P'$, where the added ghost code restricts
the path condition to heaps consistent with the chosen topology; other
topologies are explored by varying $\xi$.

\begin{wrapfigure}{r}{0.62\textwidth}
\vspace{-4mm}
{\scriptsize
\begin{minted}[escapeinside=||]{C}
void constrain_heap(Node* a0, Node* a1, Node* a2, Node* a3, int xi) {
  if (xi == 0) { /* list segment of length >= 3 */
    assume(a0 != NULL);
    assume(a1 != NULL && a2 != NULL);
    assume(a0 != a1 && a1 != a2 && a0 != a2);
    assume(a0->val == |$\beta_0$|);  assume(a0->next == a1);
    assume(a1->val == |$\beta_1$|);  assume(a1->next == a2);
    assume(a2->val == |$\beta_2$|);  assume(a2->next == a3);
  }
  /* else if (xi == 1) { ... other topologies ... } */
}
\end{minted}
}
\vspace{-4mm}
\end{wrapfigure}
\begin{example}[Ghost code for the list-topology of length $\ge 3$]
Suppose we call a procedure \texttt{foo} with a fully symbolic list head $\mathsf{call}\;\texttt{foo}(\mathsf{head}).$ We insert a ghost call right before it:
\[
\mathsf{constrain\_heap}_{\texttt{foo}}(\mathsf{a0},\mathsf{a1},\mathsf{a2},\mathsf{a3},\xi);\;
\mathsf{call}\;\texttt{foo}(\mathsf{head}).
\]
Below is one branch of $\mathsf{constrain\_heap}_{\texttt{foo}}$ implementing
the topology from Example~\ref{example:topology}. The ghost code links three concrete nodes, while leaving the list tail and the values unconstrained.
\end{example}

\begin{table}[t!]
\centering
\caption{Results on the LogicBombs benchmark.
Each program contains exactly two critical paths (bomb-triggering and benign).
Left: KLEE baselines.
Right: LLM-based tools.
Tokens are reported in millions.
The total number of bombs is 53 and the total number of critical paths (CPs) is
$53 \times 2 = 106$.}
\label{tab:logic-bombs-summary}
\footnotesize
\vspace{-2mm}
\begin{tabular}{@{}p{0.28\columnwidth}
@{\hspace{0.00\columnwidth}}
p{0.70\columnwidth}@{}}

% ================= LEFT: KLEE =================
\raggedright
\centering\textbf{KLEE Baselines}\par\vspace{2pt}
\centering
\begin{tabular}{|l|cc|}
\toprule
\multicolumn{3}{|c|}{\rule{0pt}{3.2ex}} \\  % empty header row for alignment
\textbf{Solver}
& \makecell{\textbf{Bombs}\\\textbf{Trig.}}
& \makecell{\textbf{CPs}\\\textbf{Cov.}} \\
\midrule
STP        & 20 & 59 \\
Z3         & 17 & 56 \\
Float-Z3   & 21 & 60 \\
\bottomrule
\end{tabular}

&
% ================= RIGHT: LLM TOOLS =================
\centering\textbf{LLM-Based Tools}\par\vspace{2pt}
\centering
\begin{tabular}{|l|ccc|ccc|}
\toprule
\multicolumn{1}{|c|}{\textbf{Backend}}
& \multicolumn{3}{c|}{\textbf{\projName{}}}
& \multicolumn{3}{c|}{\textbf{ConcoLLMic}} \\
\cmidrule(lr){2-4}\cmidrule(lr){5-7}
& \makecell{\textbf{Bombs}\\\textbf{Trig.}}
& \makecell{\textbf{CPs}\\\textbf{Cov.}}
& \makecell{\textbf{Tokens}}
& \makecell{\textbf{Bombs}\\\textbf{Trig.}}
& \makecell{\textbf{CPs}\\\textbf{Cov.}}
& \makecell{\textbf{Tokens}} \\
\midrule
Claude 3.7
& \textbf{50} & \textbf{103} & 0.30
& 35 & 83 & 4.03 \\

GPT-5
& 49 & 102 & 0.32
& 33 & 81 & 2.14 \\

DeepSeek-Chat
& 45 & 98 & 0.33
& 34 & 82 & 3.04 \\
\bottomrule
\end{tabular}

\end{tabular}
\vspace{-4mm}
\end{table}

\rev{Semantic heap partitioning is a heuristic: it biases exploration toward a finite set of meaningful topologies and may miss behaviors whose required topology is not generated or explored within the budget. This can yield false negatives but not false positives.}

%% \section{Part 1: LLM-Guided Inversion of Hard-to-Solve code with Iterative Bidirectional Propagation}


%% Symbolic execution often fails when encountering program fragments whose semantics are easy to compute forward but difficult for SMT solvers to solve symbolically, such as arithmetic guards, obfuscated logic, or non-linear transformations.
%% We refer to such fragments as \emph{logic bombs}~\cite{xu2018benchmarking}.
%% Suppose a program contains a hard-to-analyze statement:
%% \[
%% z = f(x),
%% \]
%% where \(f\) encapsulates complex semantics that the solver cannot symbolically reason about. So the solver is unable to find an x so that the z gets a suitable value, in a way that x and z satisfy the remaining constraints.
%% During symbolic execution, the path constraint around this function can be expressed as:
%% \[
%% \pi = \pi_{\text{prefix}} \ \land\ (z = f(x)) \ \land\ \pi_{\text{suffix}},
%% \]
%% where \( \pi_{\text{prefix}} \) collects constraints before the call to \( f \), and \( \pi_{\text{suffix}} \) those after it.  
%% When the suffix condition depends on \(z\), progress halts because the solver cannot propagate constraints backward through \(f\).

%% Our approach, \projName, addresses this by \emph{temporarily removing} the hard fragment from the symbolic path, making the affected objects symbolic, and letting the LLM synthesize a lightweight inverse or surrogate for the hard code fragment. 
%% We first replace the opaque computation with a symbolic variable:
%% \[
%% z = f(x)\;\; \leadsto \;\; \texttt{z = make\_symbolic()},
%% \]
%% allowing symbolic execution to explore the suffix freely and collect feasible concrete outputs \(z^{\ast}\) that drive the desired control flow (e.g., reaching an error branch).
%%  In this case we discover what is the value of \(z\) that satisfies the post constraints.

%% Once a target value \(z^{\ast}\) is discovered, the LLM generates a \emph{ghost inverse} \(f^{-1}\) that approximates the backward semantics of \(f\):
%% \[
%% x^{\ast} = f^{-1}(z^{\ast}).
%% \]
%% This inversion phase concretizes the inputs directly before the fragment, that would yield the observed suffix behavior, providing a bridge between solver-intractable computations and symbolic reasoning. Sou in our example we compute the value of  \(x\) that leads to our desired value for \(z\).

%% We then \emph{augment the prefix path condition}.  
%% Let $\Phi_{\mathrm{pref}}$ be the prefix constraint up to the cut-point.  
%% From the ghost execution, we know that $x$ must equal $x^{\ast}$, so we extend:
%% \[
%% \Phi_{\mathrm{prefix}}^{+}(x) \;=\; \Phi_{\mathrm{prefix}}(x)\;\wedge\;\big(x = x^{\ast}\big).
%% \]
%% Solving $\Phi_{\mathrm{prefix}}^{+}(x)$ yields concrete inputs that reproduce the same intermediate state in the original program, forcing the execution to follow the same suffix branch observed during ghost inversion.  
%% This ensures hat we define as \emph{Post-Translation Branch Suffix Equivalence (PTBSE)}; the ghost and original executions share the same suffix behavior (same branches satisfied).

%% However, the combined constraint may become unsatisfiable.  
%% Let $\Phi_{\mathrm{prefix}}(x)$ be the prefix path condition and $x = x^{\ast}$ the equality derived from inversion.  
%% If
%% \[
%% \neg\exists x.\ \Phi_{\mathrm{prefix}}(x)\ \wedge\ (x = x^{\ast}),
%% \]
%% then the inferred pre-image is unreachable under the original prefix.  
%% In that case, \projName detects UNSAT and invokes a repair procedure to find a nearby feasible value $x'$ such that
%% \[
%% \exists x.\ \Phi_{\mathrm{prefix}}(x)\ \wedge\ (x=x'),
%% \]
%% while minimizing the deviation $|x'-x^{\ast}|$.  
%% This iterative refinement process, which propagates candidate values between the prefix and suffix to make them mutually consistent, strikes a balance between satisfiability and behavioral similarity—adjusting the inverted values just enough to make the path realizable without diverging from the intended behavior.

%% From the feasible $x'$, we then compute $z' = f(x')$ and check whether $z'$ satisfies the target suffix condition.  
%% If not, we adjust and repeat.

%% The overall process forms a bidirectional refinement loop between ghost inversion and symbolic re-evaluation:
%% \begin{enumerate}
%%   \item From the suffix, collect the desired output \(z^{\ast}\) and invert to obtain a candidate input \(x^{(1)} = f^{-1}(z^{\ast})\).
%%   \item If \(x^{(1)}\) violates the prefix constraints, repair it via a distance-minimizing constraint solve to obtain \(x^{(2)}\).
%%   \item Propagate forward through \(f\) to compute \(z^{(2)} = f(x^{(2)})\); if it fails the suffix condition, treat \(z^{(2)}\) as a new soft target and invert again to produce \(x^{(3)}\).
%%   \item Continue until a consistent pair \((x^{(t)}, z^{(t)})\) satisfies both prefix and suffix constraints.
%% \end{enumerate}

%% This bidirectional reasoning—LLM-guided inversion combined with the refinement process avoids blind prefix guessing and allows symbolic execution to navigate code fragments previously opaque to solvers.
%% Rather than treating $f$ as an unmodelled black box, \projName uses ghost inversion to expose a symbolic interface for it, extending reachability to complex, solver-resistant program fragments.


%% \medskip
%% \noindent\textbf{Case 1: Simple Inversion.}  

%% \begin{lstlisting}[language=C,basicstyle=\ttfamily\small]
%% int hard_func(int x) { 
%%     return (x * x + 3) % 97; 
%% }

%% void main() {
%%     int a = input();
%%     int b = hard_func(a);
%%     if (b == 42) {
%%         assert(a > 20);
%%     }
%% }
%% \end{lstlisting}


%% During normal symbolic execution:
%% \[
%% \pi_{\text{prefix}}: \quad a = \alpha_0
%% \]
%% \[
%% y = f(a) \quad \text{where} \quad f(x) = (x^2 + 3) \bmod 97
%% \]
%% \[
%% \pi_{\text{suffix}}: \quad y = 42 \ \land\ a < 20
%% \]


%% A solver may struggle with \( y = f(a) \) since it is non-linear modulo arithmetic. 
%% Instead, the LLM-generated ghost code supplies an approximate inverse:

%% \begin{lstlisting}[language=C,basicstyle=\ttfamily\small]
%% int invert_hard_func(int y, int *results) {
%%     int count = 0;
%%     int t = (y - 3) % 97;
%%     if (t < 0) t += 97;

%%     for (int x = 0; x < 97; x++) {
%%         if ((x * x) % 97 == t) {
%%             results[count++] = x;
%%             if (count == 2) break; // mod prime => max 2 solutions
%%         }
%%     }
%%     return count;
%% }
%% \end{lstlisting}


%% We make b symbolic instead of it being calculated by the hard function and get b=42.
%% Evaluating \( f^{-1}(42) \) yields candidates \(\{10, 87\}\). The suffix constraint (\(a >20\)) prunes this to \(a = 87\), satisfying the path without heavy solver reasoning.

%% \medskip
%% \noindent\textbf{Case 2: Prefix–Suffix Conflict.}  

%% Now consider a harder variant:

%% \begin{lstlisting}[language=C,basicstyle=\ttfamily\small]
%% int a = input();
%% if (a >= 90) {
%%     int b = hard_func(a);
%%     if (b >= 42) {
%%         ...
%%     }
%% }
%% \end{lstlisting}


%% Resymbolizing $b$ yields candidate $b^* = 42$. Inversion again gives $a^* \in \{10, 87\}$.  
%% However, neither candidate is consistent: $a=10$ violates the prefix ($a > 90$), while $a=87$ also violates it. Direct inversion fails.

%% Here \projName activates our repair process. Constraints are split into:
%% \[
%% \text{Hard: } a >= 90, \qquad \text{Soft: } a = 87.
%% \]
 
%% Among assignments that satisfy the hard constraints, 
%% we choose one that minimizes the deviation from the soft constraint. 
%% In this case, the nearest feasible assignment is $a = 90$, 
%% which respects $a \geq 90$ while remaining as close as possible to $a = 87$.
%% \[
%% b = f(90) = (90^2 + 3) \bmod 97 = 52,
%% \]
%% which satisfies $b \geq 42$. The suffix is preserved, and the overall path remains feasible. 

%% If this repair had failed (e.g., yielding a $b$ that violated the suffix), \projName would iterate: softening the suffix, rerunning symbolic execution, re-inverting, and refining the candidate until both prefix and suffix converge. This loop ensures that even when inverses are approximate or conflicting, the system recovers feasible paths by balancing prefix, fragment, and suffix constraints.

%% \section{Part 2: Rewriting into Simplified Models.}
%% When inversion is infeasible, such as when multiple variables interact in non-invertible ways, or the fragment depends on external system state—\projName rewrites the code into a simplified surrogate that remains tractable for symbolic execution. 
%% The goal is not to emulate full functional fidelity, but to preserve control-relevant behavior: branches, guards, and invariants that determine the path of interest. 
%% This strategy is particularly suited to \emph{environment and system modeling} (e.g., I/O, randomness, timing, or opaque library calls), but applies equally to other solver-hostile computations.

%% \paragraph{Theoretical foundation.}
%% Rewriting relies on the notion of \emph{suffix equivalence}.  
%% Let $\pi^{\mathsf{suf}}_{f}(x)$ denote the suffix trace of program~$f$ under input~$x$, i.e., the sequence of branches or lines triggered \emph{after} a designated difficult fragment.
%% Given a rewritten program~$f'$ in which that fragment is replaced by a surrogate model, symbolic execution discovers inputs~$x'$ that reach a branch of interest.  
%% Ideally, the same input induces the same suffix behaviour in both versions:
%% \[
%% \pi^{\mathsf{suf}}_{f'}(x') \;=\; \pi^{\mathsf{suf}}_{f}(x').
%% \]
%% In practice, this equality often holds only approximately, and a remapping function~$M$ aligns surrogate inputs with feasible originals:
%% \[
%% \pi^{\mathsf{suf}}_{f'}(x') \;=\; \pi^{\mathsf{suf}}_{f}\!\big(M(x')\big),
%% \]
%% where~$M$ is derived from the surrogate and captures how the simplified computation corresponds to the real one.


%% \medskip
%% \textbf{Example 1 – File-I/O surrogate model.}
%% Consider a logging routine that writes a message to a temporary file and later reopens it for verification:
%% \begin{lstlisting}[language=C,basicstyle=\ttfamily\small]
%% write(fd, msg, strlen(msg));
%% close(fd);
%% fd = open(tmp_path, O_RDONLY);
%% read(fd, buf, sizeof(buf));
%% if (strcmp(buf, "OK") == 0)
%%     trigger_alert();
%% \end{lstlisting}

%% Here, the symbolic executor must reason about filesystem state across system calls, 
%% whether the contents written are the same bytes later read, which most symbolic engines
%% either concretize or model unsoundly.

%% \projName replaces this entire I/O chain with a compact surrogate that preserves
%% the semantic dependency between the written and read buffers:
%% \begin{lstlisting}[language=C,basicstyle=\ttfamily\small]
%% char msg[4], buf[4];
%% klee_make_symbolic(msg, 4, "msg");
%% strcpy(buf, msg);          // surrogate for write ->read round-trip
%% if (strcmp(buf, "OK") == 0)
%%     trigger_alert();
%% \end{lstlisting}

%% This rewrite eliminates all system calls yet maintains the key behavioural relation:
%% the branch depends solely on the equivalence of written and read data.  
%% The remapping function~$M$ is identity—the symbolic variable \texttt{msg}
%% represents the concrete content that would have been written and retrieved.  
%% Thus, the surrogate achieves suffix-equivalence on the branch of interest while
%% bypassing the kernel and filesystem layers entirely.

%% \medskip

%% \noindent\textbf{Example 2 – Timing surrogate.}  
%% Many programs take different paths depending on timestamps or delays:
%% \begin{lstlisting}[language=C,basicstyle=\ttfamily\small]
%% struct timeval tv;
%% gettimeofday(&tv, NULL);
%% if (tv.tv_sec % 2 == 0) do_even(); else do_odd();
%% \end{lstlisting}

%% \projName abstracts the system clock as a bounded symbolic integer:
%% \begin{lstlisting}[language=C,basicstyle=\ttfamily\small]
%% int sec;
%% klee_make_symbolic(&sec, sizeof(int), "sec");
%% klee_assume(sec >= 0 && sec < 60);
%% if (sec % 2 == 0) do_even(); else do_odd();
%% \end{lstlisting}

%% The surrogate eliminates the external time source while maintaining
%% suffix-equivalence on both branches.

%% \medskip
%% \noindent\textbf{Example 3.} 
%% This approach extends beyond environment modeling: it can be applied to transform solver-hostile code fragments into solver-friendly surrogates that mimic the original functionality—or at least preserve a semantically relevant subset of it.
%% Consider a floating-point computation that SMT solvers cannot easily handle:
%% \[
%% y = x \cdot 1.035 , 
%% \quad 
%% \texttt{if (y > 400) \{ ... \}}.
%% \]
%% We introduce a surrogate using integer arithmetic:
%% \[
%% y = x \cdot 1035.
%% \]
%% Symbolic execution on $f'$ may find an input $x'$ such that $y > 50$ in 
%% the surrogate. However, this $x'$ will not directly satisfy the floating-point 
%% condition in the original program (for example $x'=4$ ). Instead, a remapping is required, for example
%% \[
%% M(x') = 100 \cdot x',
%% \]
%% which realigns the surrogate inputs with the original semantics. 

%% Thus, suffix equivalence is preserved up to $M$, ensuring that critical 
%% branches covered in the surrogate program correspond to valid branches 
%% in the original.

%% Fora second example, c
%% Thus, the surrogate preserves the triggering of the critical branch without introducing 
%% spurious executions. Symbolic execution can then proceed on the surrogate program, free of 
%% transcendental constraints, while \emph{Post-Translation Path Suffix Equivalence (PTPSE)} ensures 
%% that for inputs that reach the branch of interest in the simplified model , we are going to trigger the same branch in the original program.

%% \section{Part 3: Dynamic Memory and Function pointers}
%% \label{sec:pointers}
%% \noindent
%% This part develops a single idea in two settings: \emph{bound the unknown space with a small, semantically meaningful menu and expose it as a symbolic choice}. First, for dynamic memory, we ask the LLM to partition heap behaviors into a few canonical \emph{topologies} that matter for control (e.g., ring vs.\ chain, below vs.\ above a limit). Each topology is materialized by a compact ghost builder and guarded by a bounded symbolic chooser, so the executor ranges over a finite set of representative shapes instead of enumerating arbitrary aliasings. Second, for indirect calls (function pointers and virtual dispatch), we apply the same recipe: a lightweight analysis gathers a sound over-approximation of potential callees, the LLM ranks it to a small \emph{candidate set}, and a single symbolic index selects the active target. In both cases, the ghost layer preserves program semantics, keeps interior details symbolic, and makes exploration \emph{bounded, replayable, and refinable}, turning intractable heap search and opaque dispatch into tractable, semantically curated case splits.


%% \subsection{Dynamic Memory Topological partition for Symbolic Execution}
%% \label{sec:ghost-guided-heap}
%% Symbolic execution performs well when control flow is governed by scalar inputs, but its effectiveness deteriorates once it depends on heap wiring—aliasing relationships (who aliases whom, which node a pointer refers to) or rare memory configurations that unlock deep paths. In practice, when dynamic memory is involved, three factors commonly hinder coverage: 
%% \begin{itemize}
%%     \item \textbf{Path explosion} from symbolic pointers, where each dereference multiplies states by the number of plausible pointees.
%%     \item \textbf{Solver strain} from pointer-encoded constraints, where nested memory-array selects are difficult to invert.
%%     \item \textbf{Lack of data-structure semantics}, which leads engines to waste budget exploring isomorphic heaps, cycles, and back-edges that do not affect control.
%% \end{itemize}
%% As a result, pointer-centric code often hides witnesses behind complex heap shapes and aliasing relations that the sym-ex engine is unlikely to assemble within realistic resource limits.

%% Our approach does not aim to replace the solver with an LLM, but rather to provide the solver with a smaller, more meaningful search space. In essence, the LLM is used at the beginning of the process, to impose a semantically bounded abstraction of heap behaviors. Specifically, the model synthesizes a lightweight “ghost” layer composed of: (a) compact builders that materialize a few canonical heap patterns relevant to control flow, and (b) a small number of finite-domain \emph{symbolic choosers}, which are bounded symbolic integers that select among those patterns or precomputed hit/miss outcomes for solver-hostile guards. All other fields remain symbolic. This transforms the problem from “guessing addresses until the heap happens to align” into “selecting among a few semantically distinct cases and letting the solver handle the remaining constraints.” The original program logic is preserved, witnesses are replayable, and the number of additional forks is bounded by construction, while symbolic execution continues to perform feasibility checking within each selected class.


%% \subsubsection{Formalization: Dynamic Memory Topological Partition}
%% \label{subsubsec:formalization}

%% We begin with an imperative program $P$ and a some coverage data from baseline symbolic execution runs. 
%% Let $S$ denote the set of \emph{control-relevant sites}; pointers, heap structures, or fields that influence branching, and let $F$ denote the set of \emph{solver-hostile guards}, such as non-linear arithmetic or nested array selects. 
%% These are identified through coverage information: branches that remain uncovered or guards that are repeatedly pruned or timed out during baseline symbolic execution. 
%% They guide where the LLM should focus its abstraction effort, prioritizing difficult paths to improve coverage and resource allocation.

%% \paragraph{LLM-guided semantic partitioning.}
%% Given a target function or module under test, we provide the LLM with its body together with the relevant context required to understand its control flow.  
%% This context may, in principle, be assembled using sophisticated retrieval or dependency-analysis techniques; in our prototype, we adopt a simple yet effective strategy: we include all global variables and data structures that the function accesses, along with the first-order callees it directly invokes.  
%% This keeps each prompt compact and well-scoped, enabling the process to remain tractable and scalable across large codebases while still exposing sufficient semantic information for the model to infer control-relevant heap behaviours.

%% Using this information, , the LLM performs a one-shot semantic analysis of how
%% dynamic memory configurations influence branching.
%% For example, in a queueing function, distinct cases may involve a buffer below or above a limit; in a list-manipulation function, they may involve acyclic lists, cyclic lists, or lists exceeding a threshold length.  
%% This results in a finite \emph{semantic topological partition} of the heap space:
%% \[
%% \Pi = \{\, C_1, C_2, \dots, C_k \,\},
%% \]
%% where each class $C_i$ captures a family of heap and input configurations that yield a distinct control-flow behaviour.  
%% Two classes are considered distinct if they induce different outcomes for at least one branch whose condition depends on heap topology or pointer relations.  
%% The LLM thus implements a function
%% \[
%% \mathsf{LLM\_partition}(P, S, F) \rightarrow \Pi,
%% \]
%% where $S$ is the set of control-relevant sites—pointers, heap structures, or fields that influence branching—and $F$ is the set of solver-hostile guards, such as non-linear arithmetic or nested array selects, as defined earlier.
%% These two sets capture, respectively, the logic of heap-dependent control and the expressions that are difficult for constraint solvers to invert, guiding the model’s semantic partitioning of the search space into a finite set of representative classes~$\Pi$.


%% \paragraph{From partitions to ghost builders.}
%% Each class $C_i$ is then translated into a small executable fragment, called a \emph{ghost builder}, that constructs a representative heap shape satisfying the intended properties of that class.  
%% Formally, each builder is associated with a predicate
%% \[
%% \varphi_i(H,\nu),
%% \]
%% where $H$ denotes the heap—encoding the topology and contents of dynamically allocated objects, and $\nu$ denotes the environment mapping program variables to their current (symbolic) values.  
%% The predicate $\varphi_i$ characterizes the conditions under which class~$C_i$ holds, potentially involving both the heap (e.g., “the list rooted at \texttt{p} is cyclic”) and scalar program state (e.g., “the counter variable \texttt{count} exceeds its limit”).  
%% The builder $\mathsf{build}_i()$ produces a heap~$H$ such that $\varphi_i(H,\nu)$ is satisfied, while keeping all other values symbolic.
%% Importantly, builders are not test cases: they do not assign concrete data values, but only enforce the minimal structural invariants that distinguish one class from another. 
%% % For example the length of the list will be set as a symbolic value l, which satisfies the property l > desired length.
%% All remaining quantities, such as payloads, scalar values, or irrelevant fields, remain symbolic, allowing the solver to explore within each class freely. 

%% To mitigate solver strain from pointer-encoded constraints (e.g., nested selects or difficult inversions), builders may also emit small “inversion hints,” precomputing indices or satisfying assignments for specific conditions that would otherwise require array reasoning.  
%% This transformation converts solver-hostile formulas into lightweight, executable recipes.

%% A bounded symbolic integer $\chi \in \{1,\dots,k\}$, called the \emph{chooser}, determines which builder is active at runtime.  
%% The resulting \emph{ghost-augmented program} is:
%% \[
%% G(P): \quad \texttt{switch}(\chi)\ \{\ 
%%     \mathsf{build}_1();\ \mathsf{build}_2();\ \dots\ \mathsf{build}_k();\ \},
%% \]
%% where each branch materializes one canonical heap topology.  
%% Symbolic execution then ranges over the bounded domain of~$\chi$, exploring each class exhaustively.  
%% Each configuration remains symbolic internally, e.g., a list length $s>N$ is enforced symbolically via \texttt{klee\_assume($s > N$)}, so coverage within each class is preserved.

%% When $k$ is large, ghost builders can be generated incrementally: each new class simply adds a case guarded by an additional symbolic value, extending the bound of~$\chi$.  
%% Because builders are compositional and independent, the synthesis remains scalable even across multiple LLM invocations.  
%% The resulting ghost layer preserves the semantics of the original program while making the exploration space explicitly bounded and incrementally extensible.  
%% Each ghost builder is concretizable: any satisfying model discovered by symbolic execution on~$G(P)$ can be re-executed on~$P$ to reproduce the same control-flow behavior, ensuring \emph{replayability}.  
%% Ghost code is inserted only at heap-materialization points and never alters the program’s functional logic, maintaining \emph{non-interference} and observational equivalence with~$P$ under any fixed chooser valuation.  
%% Because each semantic class and its symbolic chooser are explicit, the partition can be refined locally: if coverage feedback indicates an unrepresented behavior, a new class can be added without disturbing existing ones, providing natural \emph{refinability}.  
%% This transformation balances expressiveness and practicality: if a real execution of~$P$ falls within a represented semantic class $C_i \in \Pi$, symbolic execution of~$G(P)$ can reach it, while the number of additional symbolic forks grows only linearly with~$k$, independent of heap size or aliasing degree.  
%% The LLM’s role is limited to synthesizing a finite semantic partition and lightweight builders; it does not perform constraint solving or generate concrete tests—symbolic execution and the solver remain responsible for exploring feasibility within each class.  
%% While our prototype uses KLEE, the approach applies to any symbolic executor that forks on symbolic pointers or models heap memory through arrays.  
%% Overall, the method offers a systematic bridge between program semantics and the bounded search capabilities of symbolic execution, yielding a semantically grounded yet computationally tractable exploration process.

%% %% \subsubsection{Motivating Example 1: Memory-Dependent Branch in a Networking Queue}

%% %% A practical case of heap-dependent control flow appears in the enqueue fast path of Linux’s \texttt{sch\_fq\_codel} discipline, where packets are either accepted or dropped depending on both queue length and total buffered bytes. The core logic can be distilled to:
%% %% \begin{lstlisting}[language=C, basicstyle=\ttfamily\small]
%% %% int fq_codel_enqueue(struct skb *skb, struct Qdisc *sch,
%% %%                      struct fq_codel_sched_data *q) {
%% %%     q->memory_usage += skb->truesize;       // (1) total bytes
%% %%     if (++sch->q.qlen <= sch->limit &&
%% %%         q->memory_usage <= q->memory_limit)
%% %%         return NET_XMIT_SUCCESS;            // accept
%% %%     drop_one(sch, q);
%% %%     if (q->memory_usage > q->memory_limit)
%% %%         q->drop_overmemory++;               // rare branch
%% %%     return NET_XMIT_CN;
%% %% }
%% %% \end{lstlisting}

%% %% The rarely taken branch at the end triggers only when the aggregate size of all queued packets exceeds the configured memory limit while remaining below the packet-count cap.  
%% %% A conventional KLEE harness would construct a symbolic list of packets with symbolic \texttt{truesize} fields (1–32\,KB) and attempt to discover a combination whose total sum crosses 32\,MB. Even for moderate queue lengths, this produces an exponential constraint space: the solver wastes time exploring countless short queues before ever exceeding the threshold, a pattern common in memory-guarded logic.

%% %% In the ghost-guided variant, the symbolic list is replaced by a compact builder that materializes the aggregate usage directly:
%% %% \begin{lstlisting}[language=C, basicstyle=\ttfamily\small]
%% %% static void fill_usage(unsigned bytes){
%% %%     while (bytes > 0) { g_q.memory_usage += 32*1024; bytes -= 32*1024; g_sch.q.qlen++; }
%% %% }
%% %% \end{lstlisting}
%% %% A single symbolic chooser then partitions execution into two semantically distinct classes:
%% %% \begin{lstlisting}[language=C, basicstyle=\ttfamily\small]
%% %% if (choose == 0) klee_assume(usage < MEM_LIMIT / 2);          // under limit
%% %% else              klee_assume(usage > MEM_LIMIT - 32*1024*10); // near limit
%% %% fill_usage(usage);
%% %% \end{lstlisting}

%% %% Instead of exploring thousands of unconstrained heap configurations, symbolic execution now examines only two meaningful cases, normal operation and near-limit exhaustion, covering both behaviors instantly while preserving the original semantics.  
%% %% This example illustrates how bounded symbolic choosers and lightweight heap builders can turn intractable, aggregate heap conditions into a small number of semantically representative configurations.


%% %% \subsubsection{Motivating Example 2: Topology-Exclusive Feasibility — A Two-Level Skip-List Ring}

%% %% Consider a two-level ``skip-list–like'' structure in which each node has a level-0 next pointer (\texttt{lvl0}) and a level-1 express pointer (\texttt{lvl1}). They also have a integer value field, and a char array name field.  
%% %% The function under test, \texttt{skiplist\_hard(h, g)}, verifies purely topological properties: (1)~the \texttt{lvl0} pointers form a 4-cycle; (2)~each \texttt{lvl1} points two hops ahead; and (3)~following two \texttt{lvl0} links from $g$ returns to $h$.  
%% %% Any violation returns~0; the success case issues an assertion.

%% %% \begin{lstlisting}[language=C,basicstyle=\ttfamily\small]
%% %% int skiplist_hard(SNode *h, SNode *g) {
%% %%   SNode *a=h?h->lvl0:NULL, *b=a?a->lvl0:NULL;
%% %%   SNode *c=b?b->lvl0:NULL, *d=c?c->lvl0:NULL;
%% %%   if (!h||!a||!b||!c||!d) return 0;
%% %%   if (d!=h) return 0;
%% %%   if (h->lvl1!=b||a->lvl1!=c||
%% %%       b->lvl1!=d||c->lvl1!=a) return 0;
%% %%   SNode *g1=g?g->lvl0:NULL,*g2=g1?g1->lvl0:NULL;
%% %%   if (g2!=h) return 0;
%% %%   klee_assert(0 && "ring+express+meet");
%% %% }
%% %% \end{lstlisting}

%% %% A naïve symbolic harness allocates a pool of nodes and makes every pointer field symbolic.  
%% %% The solver must then satisfy the ring equalities (\texttt{lvl0}$^4$=h), four express-alignment constraints (\texttt{lvl1}= \texttt{lvl0}$^2$ for each node), and the cross-head condition (\texttt{g\_lvl0\_lvl0}=h).  
%% %% These form a narrow manifold in heap space—most random pointer combinations yield \texttt{NULL}s or broken chains, so the search budget  can be exhausted long before assembling the exact configuration that satisfies all relations.

%% %% The ghost-guided version replaces the open-ended heap search with a single bounded symbolic chooser $\texttt{klass}\in\{1,\dots,5\}$ that selects among five canonical topologies: broken chain, non-ring, misaligned express pointers, mismatched head link, and full success.  
%% %% Within each class, only the relevant links are fixed; all other fields remain symbolic, including the value and name fields.  
%% %% Symbolic execution thus explores five semantically distinct cases instead of an exponential pointer product, reliably reaching the success path.

%% %% Figure~\ref{fig:skiplist-ring} shows the target success topology (Case~5): a 4-cycle on \texttt{lvl0}, express pointers two hops ahead, and $g$ placed so that two \texttt{lvl0} steps land at $h$.
%% %% \vspace{0.5em}
%% %% \begin{figure}[h]
%% %%   \centering
%% %%   \resizebox{1.16\linewidth}{!}{%
%% %% \begin{tikzpicture}[
%% %%     >=stealth,
%% %%     node distance=2cm,
%% %%     main/.style={draw,rounded corners,minimum width=11mm,minimum height=8mm,align=center},
%% %%     arr/.style={->,very thick},
%% %%     exp/.style={->,densely dashed,thick,black!70},
%% %%     gptr/.style={->,ultra thick,red!70!black,shorten >=1pt,shorten <=2pt},
%% %%     x=1cm,y=1cm
%% %%   ]
%% %%   \centering
%% %%   % --- main nodes ---
%% %%   \node[main] (n0) at (-3, 1.5) {n0};
%% %%   \node[main] (n1) at ( 3, 1.5) {n1};
%% %%   \node[main] (n2) at ( 3,-1.5) {n2};
%% %%   \node[main] (n3) at (-3,-1.5) {n3};

%% %%   % --- lvl0 ring: solid inner square ---
%% %%   \draw[arr] (n0) -- node[draw=none,fill=none,above]{lvl0} (n1);
%% %%   \draw[arr] (n1) -- node[draw=none,fill=none,right]{lvl0} (n2);
%% %%   \draw[arr] (n2) -- node[draw=none,fill=none,below]{lvl0} (n3);
%% %%   \draw[arr] (n3) -- node[draw=none,fill=none,left]{lvl0} (n0);

%% %%     \node[draw=none,fill=none,above left=15pt and 7pt of n0] (glabel) {$h$};
%% %%   \draw[gptr] (glabel.north east) -- (n0.north);

%% %%   % --- lvl1 express links: dashed, routed with outer control points ---
%% %%   \draw[exp]
%% %%     (n0) .. controls ($(n0)+(0,3.8)$) and ($(n2)+(0,3.8)$) ..
%% %%     node[draw=none,fill=none,above]{lvl1} (n2.north west);
%% %%   \draw[exp]
%% %%     (n1) .. controls ($(n1)+(5.8,0)$) and ($(n3)+(5.8,0)$) ..
%% %%     node[draw=none,fill=none,right]{lvl1} (n3.north east);
%% %%   \draw[exp]
%% %%     (n2) .. controls ($(n2)+(0,-3.8)$) and ($(n0)+(0,-3.8)$) ..
%% %%     node[draw=none,fill=none,below]{lvl1} (n0.south east);
%% %% \draw[exp]
%% %%     (n3) .. controls ($(n3)+(-5.8,0)$) and ($(n1)+(-5.8,0)$) ..
%% %%     node[draw=none,fill=none,left]{lvl1} (n1.south west);



%% %%   % --- g pointer (label only, no box) ---
%% %%   \node[draw=none,fill=none,below right=15pt and 7pt of n2] (glabel) {$g$};
%% %%   \draw[gptr] (glabel.north west) -- (n2.south);
%% %% \end{tikzpicture}
%% %% }
%% %% \caption{Solid edges (\texttt{lvl0}) form a 4-cycle; dashed outer edges (\texttt{lvl1}) skip two nodes ahead; the red arrow shows pointer \texttt{g} to \texttt{n2}.}
%% %% \label{fig:skiplist-ring}
%% %% \end{figure}

%% %% Naïve symbolic execution wastes effort enumerating millions of invalid aliasings before ever assembling the ring+express+meet configuration.  
%% %% The ghost-guided harness constrains exploration to five semantically distinct topologies—a small, meaningful partition that makes the success path routinely reachable while leaving non-topological fields symbolic.


%% %% \subsubsection{Motivating Example 3: Topology-Dependent Merge in an \texttt{ext4} Extent Tree}

%% %% Our final example comes from Linux’s \texttt{ext4\_ext\_insert\_extent()} routine, which inserts a new \emph{extent}—,  contiguous sequence of physical disk blocks, into a leaf of the extent tree.  
%% %% Each leaf describes ranges of logical file blocks and their corresponding physical locations on disk.  
%% %% When two consecutive extents occupy adjacent logical blocks and point to consecutive physical blocks, they can be merged to reduce fragmentation and metadata overhead.

%% %% \begin{lstlisting}[language=C,basicstyle=\ttfamily\small]
%% %% int ext4_ext_insert_extent(struct ext4_ext_path *path,
%% %%                            const struct ext4_extent *newext_in){
%% %%   struct ext4_extent_header *eh = path->p_hdr;
%% %%   struct ext4_extent *ex = path->p_ext;
%% %%   if (!eh || le16_to_cpu(eh->eh_magic)!=EXT4_EXT_MAGIC) return -1;
%% %%   if (le16_to_cpu(eh->eh_depth)!=0) return -2;
%% %%   int entries = le16_to_cpu(eh->eh_entries);
%% %%   if (entries<=0) return -3;
%% %%   int idx = (int)(ex - (struct ext4_extent*)(eh+1));
%% %%   if (idx<0 || idx>=entries) return -4;
%% %%   struct ext4_extent newext=*newext_in;
%% %%   if (ext4_can_extents_be_merged(ex,&newext)){
%% %%     ex->ee_len = cpu_to_le16(le16_to_cpu(ex->ee_len)
%% %%                             + le16_to_cpu(newext.ee_len));
%% %%     ext4_ext_try_to_merge_right(path);    // optional compaction
%% %%     return 0;                             // merged
%% %%   }
%% %%   return 1;                               // no merge
%% %% }
%% %% \end{lstlisting}

%% %% Here, \texttt{ext4\_can\_extents\_be\_merged(a,b)} checks whether the two extents are contiguous both in logical block space and on disk.  
%% %% \emph{Logical adjacency} means that the new extent’s starting logical block immediately follows the previous extent’s range, while \emph{physical adjacency} means that their corresponding disk blocks are consecutive.  
%% %% Only when both hold can the extents be safely merged and optionally compacted by \texttt{ext4\_ext\_try\_to\_merge\_right()}.

%% %% When all fields and pointers are symbolic, these cross-object equalities define a narrow feasible region: most models violate one of the adjacency conditions or early structural guards.  
%% %% As a result, symbolic execution quickly reaches the early-return and no-merge cases but rarely discovers the merge or compaction paths within realistic resource limits.


%% %% \paragraph{Generic LLM Prompt Template.}
%% %% To obtain the semantic topological partition automatically, \projName
%% %% invokes the LLM with a generic instruction that is agnostic to the specific
%% %% function being analyzed. The only input is the code fragment itself and
%% %% (optionally) short contextual comments or coverage statistics.
%% %% The model is asked to describe distinct heap or pointer configurations
%% %% that could lead to different control-flow outcomes.

%% %% \begin{lstlisting}[language=text,basicstyle=\ttfamily\small]
%% %% Task: Identify a small set of semantically distinct heap or pointer configurations ("classes") that capture different control-flow behaviors in the given function. Each class should represent a canonical topology
%% %% (e.g., valid, invalid, boundary, merged, cyclic, empty) or a distinct semantic case (e.g., success/failure, under/over  threshold). Use concise natural-language bullets.
%% %% Input: [insert code of the target function here]

%% %% Goal:
%% %% - Produce k minimal, non-overlapping classes C1..Ck
%% %% - Each class should define what heap or pointer relations hold   (e.g., which pointers alias, which structures are null or cyclic)
%% %% - Keep k small but sufficient to isolate all unique control-flow outcomes.
%% %% \end{lstlisting}

%% %% Applied to the \texttt{ext4\_ext\_insert\_extent()} routine, this prompt
%% %% yields the following natural-language partition (Step~1 output), which
%% %% the LLM then later transforms into executable ghost builders (Step~2).


%% %% \paragraph{LLM NL Output (Step~1: partition)} .

%% %% \begin{lstlisting}[language=text,basicstyle=\ttfamily\small]
%% %% C1 (Invalid header): Header missing or bad magic => early return -1
%% %% C2 (Not a leaf): The eh_depth is not zero => early return -2
%% %% C3 (Empty leaf): The number of eh_entries is less or equal to zero => early return -3
%% %% C4 (Bad index): p_ext out of bounds => early return -4
%% %% C5 (Merge + right-compaction): ex and new are adjacent both logically and physically; 
%% %%     right neighbor is also adjacent so ext4_ext_try_to_merge_right folds it => merged+compacted path
%% %% C6 (Merge only): ex and new are adjacent (logical & physical), but the right neighbor is not 
%% %%     => merged path, no compaction
%% %% C7 (No merge): Break either logical or physical adjacency, while other structure is valid 
%% %%     => return 1
%% %% \end{lstlisting}


%% %% Based on these partitions the LLM then can produce the ghost version which replaces the unstructured heap search with a single bounded symbolic chooser \(\texttt{klass}\!\in\!\{1,\dots,7\}\) selecting among canonical leaf layouts:  
%% %% invalid header, non-leaf, empty leaf, bad index, merge + compaction, merge-only, and no-merge.  The minimal harness is shown below, with the make\_adjacent and make\_not\_adjacent omitted for brevity.
%% %% \paragraph{LLM produced Ghost Code (Step~2: Ghost generation)}
%% %% .
%% %% \smallskip

%% %% \begin{lstlisting}[language=C,basicstyle=\ttfamily\small]
%% %% int klass = klee_range(1, 8, "klass");
%% %% build_header(&eh); build_extents(arr);
%% %% switch (klass){
%% %%   case 1: path.p_hdr=NULL; break;
%% %%   case 2: eh->eh_depth=1; break;
%% %%   case 3: eh->eh_entries=0; break;
%% %%   case 4: path.p_ext=&arr[eh->eh_entries]; break;
%% %%   case 5: make_adjacent(arr[1], newext, true); break;   // merge+right
%% %%   case 6: make_adjacent(arr[1], newext, false); break;  // merge-only
%% %%   case 7: make_nonadjacent(arr[1], newext); break;      // no merge
%% %% }
%% %% (void)ext4_ext_insert_extent(&path, &newext);
%% %% \end{lstlisting}

%% %% Symbolic execution now explores seven bounded, semantically distinct configurations instead of unconstrained heaps, consistently reaching the merge and compaction cases.  
%% %% The ghost layer does not alter logic.it merely structures the heap search into a finite, replayable partition, making this topology-dependent path both sound and tractable.

%% \subsection{Function Pointers and Virtual Calls}

%% Indirect calls through function pointers or C++ virtual tables present obstacles to symbolic execution closely analogous to the dynamic-memory issues we studied. In practice, we observed two recurring failure modes with \textsc{KLEE} on real C/C++ code: (i) \emph{opaque callees}, where unresolved helpers or vtable slots force the engine to concretize or stub the target unsoundly; and (ii) \emph{explosive over-approximation}, where attempting to be sound by enumerating all compatible targets yields an intractable number of paths. Both patterns are common in registries, plugin APIs, and large class hierarchies.

%% \medskip
%% \noindent\textbf{Ghost-guided candidate sets.}
%% As with dynamic memory, we \emph{bound and expose} the unknown space explicitly. A lightweight analysis gathers a sound but coarse set of potential callees (matching prototype and calling convention, addresses-taken only; for virtuals, overrides of the same slot). An LLM then ranks this set semantically (names, comments, usage sites) to keep the top-$k$ plausible targets (typically $k\!\in\![3,5]$). The indirect call is rewritten so the callee is selected by a small symbolic index, yielding exactly $k$ forks.

%% \medskip
%% \noindent\textbf{Motivating example: math “strategy” callback.}
%% Consider an analytics kernel that dispatches to one of many math strategies selected at runtime (configuration, plugin, or build flag). Different strategies have different domain behaviors (NaN propagation, saturation, divide-by-zero handling), and these differences drive control flow:

%% \begin{lstlisting}[language=C, basicstyle=\ttfamily\small]
%% typedef double (*binop_t)(double, double);

%% static double add_(double x, double y){ return x + y; }
%% static double mul_(double x, double y){ return x * y; }
%% static double div_safe_(double x, double y){ return (y==0)? NAN : x / y; }
%% static double pow_sat_(double x, double y){
%%   double r = pow(x, y);
%%   return isfinite(r) ? r : copysign(DBL_MAX, r);
%% }
%% static double hypot_(double x, double y){ return hypot(x, y); }
%% ... other potential math functions

%% double score(binop_t op, double a, double b, double thr){
%%   double v = op(a, b);
%%   if (isnan(v))            return -INFINITY;   // error path
%%   if (v > thr)             return 1.0;         // accept
%%   return 0.0;                                   // reject
%% }
%% \end{lstlisting}

%% In a realistic build there may be dozens of address-taken \texttt{binop\_t} implementations (plugins, feature toggles, SIMD variants). Naïve \textsc{KLEE} either concretizes \texttt{op} (missing paths where, e.g., \texttt{div\_safe\_} yields \texttt{NaN}) or explodes when forking one path per candidate. 

%% \medskip
%% Let $\mathcal{C}(\texttt{fp}) = \{t_0, t_1, \dots, t_{k-1}\}$ denote the finite
%% \emph{candidate set} of possible callees for a function pointer~\texttt{fp},
%% selected by our static analysis and semantically ranked by the LLM.
%% A fresh bounded symbolic integer $i_{\texttt{fp}} \in \{0, \dots, k{-}1\}$
%% represents the active target chosen at runtime.
%% The indirect call is rewritten as:
%% \[
%% \texttt{ret} := \mathcal{C}(\texttt{fp})[\,i_{\texttt{fp}}\,](a_1, a_2, \ldots, a_n),
%% \]
%% where $(a_1,\dots,a_n)$ are the original call arguments and
%% $\mathcal{C}(\texttt{fp})[\,i_{\texttt{fp}}\,]$ denotes the function
%% $t_{i_{\texttt{fp}}}$ selected from the candidate set.

%% This mirrors our dynamic-memory approach: a small,
%% semantically curated candidate set guarded by a symbolic chooser turns otherwise
%% unreachable indirect-call paths into tractable, reproducible test cases without
%% altering the program’s functional logic.

%% So for the previous example, the ghost-guided rewrite bounds the dispatch to a ranked short list capturing distinct control effects (NaN propagation, saturation, geometric combination):

%% \begin{lstlisting}[language=C, basicstyle=\ttfamily\small]
%% /* Coarse set (prototype match & address-taken), ranked to top-k by the LLM */
%% static binop_t C[3] = { div_safe_, pow_sat_, hypot_ }; // k = 3

%% int idx = klee_range(0, 3, "op_idx");
%% binop_t op = C[idx];
%% (double)score(op, a, b, thr);
%% \end{lstlisting}

%% Symbolic execution now explores exactly three semantically distinct behaviors: (1) \texttt{div\_safe\_} can drive the \texttt{isnan} error branch via divide-by-zero; (2) \texttt{pow\_sat\_} reaches acceptance by saturating large exponents; and (3) \texttt{hypot\_} may remain below threshold unless inputs are large. The cost of indirect dispatch becomes a single bounded integer, not symbolic reasoning about code addresses or an unbounded fork over all plugins. As before, replayability holds: the model of \texttt{op\_idx} pins the callee, and the run reproduces on the unmodified program.



\section{Evaluation}
\label{sec:evaluation}

Our experiments address the following research questions:
\begin{description}
\item[RQ1:] Do LLM-generated ghost code techniques in \projName unlock solver-hostile paths and improve coverage compared to existing symbolic and LLM-based approaches?
\item[RQ2:] Does \projName scale to real-world numerical and structured-input programs?
\item[RQ3:] How does LLM token cost of \projName compare to prior LLM-based approaches?
\item[RQ4:] \rev{How quickly does \projName{} accumulate coverage relative to baselines?}
\item[RQ5:] \rev{Does \projName{} explore hard-to-reach paths that other techniques miss?}
\end{description}

\paragraph*{Experimental Setup} \projName is implemented on top of KLEE~\cite{klee} with Z3~\cite{de2008z3}, and KLEE-Float was used on instances involving floating point numbers. All experiments were conducted on a single Linux machine running Ubuntu 24.04.3 LTS with an Intel® Core™ Ultra 7 258V CPU and 32~GiB of RAM. All symbolic execution runs were conducted inside Docker containers to provide isolation and reproducibility. Each tool configuration was allocated a fixed wall-clock budget of 36 hours, which was applied uniformly across all tools, baselines, and ablations, following standard practice in symbolic execution and software testing evaluations.

\paragraph*{LLM Backends.} All LLM inference was performed via external APIs; we used three state-of-the-art LLMs from different families: (1) \textbf{GPT-5} (version \texttt{gpt5-2025-08-07})~\cite{openai_gpt5}, (2) \textbf{Claude-3.7} (version \texttt{claude-3-7-sonnet-20250219})~\cite{anthropicclaude}, and (3) \textbf{DeepSeek-Chat} (version 3.2, release 2025-12-01)~\cite{deepseekv3}. We used the sampling temperature of 0.5, a standard mid-range setting, and to enhance reproducibility, we cached all LLM responses so that they can be replayed deterministically~\cite{dai2025statisticalindependenceawarecaching}.

%\rev{The artifact also includes the full prompts, cached LLM responses, generated ghost code, harnesses, logs, and replay scripts.}

\paragraph*{Benchmarks} We used three benchmarks to stress different symbolic execution challenges:
\begin{itemize}
\item \textbf{LogicBombs} benchmark suite introduced by Xu et al.~\cite{xu2018benchmarking} consists of small, synthetic C programs embedding guarded ``logic bombs'', each targeting a specific symbolic execution challenge (e.g., arithmetic reasoning, floating-point conditions, symbolic memory, and path explosion). In total, it comprises 53 programs with 1{,}121 lines of executable code; each of them isolates a single challenge, enabling controlled analysis of which reasoning techniques are effective and how individual components of \textsc{\projName{}} contribute to the overall result.
\item \textbf{FDLibM} (Freely Distributable LibM)~\cite{FDLibM} is a popular math library involving complex arithmetic and floating-point guards. Our evaluation covers 78 entry points, exercising different parts of the library, totaling 5{,}013 lines of executable code. FDLibM contains deeply nested control flow, non-linear arithmetic, and floating-point branching that frequently leads to concretization or solver failure in conventional symbolic execution. This benchmark therefore allows us to evaluate \projName's ability to reason about inversion patterns, mathematical structure, and floating-point edge cases in real-world numerical code.
 \item \textbf{Structured-Input Programs}: to evaluate \projName capabilities on larger, real-world codebases, we consider programs that process structured input, since such programs rely on complex dynamic data structures. We chose \rev{four} widely-used programs: (1) \texttt{libexpat} (XML parsing, 11{,}828  LOC), (2) \texttt{jq} (JSON processing, 3{,}399  LOC), (3) GNU \texttt{bc} (POSIX expression language, 8{,}182  LOC), \begin{revblock} and (4)  \texttt{libyaml} (YAML parsing, 2{,}760  LOC).   \end{revblock}
\end{itemize}

\begin{figure}[t]
  \centering
 \includegraphics[
  width=0.9\textwidth,
  trim=0 8 0 0,
  clip
 ]{technique_queries_composite.png}
 \vspace{-3mm}
  \caption{Coverage across different combinations of tool components. 
  Dots indicate enabled techniques, while bars report the number of logic bombs triggered.
  Black rows show individual techniques, purple rows show intersections (AND),
  blue rows show unions (OR), and the green row highlights the union of the three
  LLM-based techniques, which is equivalent to \projName.}
  \label{fig:technique-queries}
 \vspace{-3mm}
\end{figure}


\paragraph*{Baselines} We compared \projName against the following baselines:
\begin{itemize}
\item Traditional symbolic execution: (1) \textbf{KLEE 3.0 (STP)}, using the official \texttt{klee:3.0} Docker image with the STP solver, (2) \textbf{KLEE 3.0 (Z3)}, using the same image with Z3 as the backend solver, and (3) \textbf{KLEE-Float (Z3)}, using a fork of KLEE \cite{kleeFloat} designed for precise floating-point reasoning. All tools used the default settings provided by the corresponding distributions, using exactly the same options when KLEE serves as the \projName backend. \rev{Thus, unless explicitly stated otherwise, \projName{} uses the same KLEE/Z3 options as the corresponding KLEE baseline.}
\item \textbf{ConcoLLMic}~\cite{concollmic} as a representative LLM-based concolic execution baseline. It first performs code instrumentation, followed by LLM-guided test generation and path exploration. \rev{ConcoLLMic's instrumentation phase can consume much or all of the time budget before any test generation begins. To objectively assess its exploration capability, we exclude ConcoLLMic's instrumentation time from the 36-hour budget---a choice that favors ConcoLLMic over \projName{}.} We evaluated it using the same three LLMs: GPT-5, Claude-3.7, and DeepSeek-Chat.
\item \textbf{Cottontail}~\cite{cottontail} is an LLM-based concolic executor designed for grammar-constrained input generation. Thus, we compared its performance with \projName on the structured-input programs (\texttt{libexpat}, \texttt{jq}, \texttt{bc}, \texttt{libyaml}) using the same LLMs. We used JSON and XML grammars provided by the original implementation; and added grammar support for \texttt{bc} and YAML.
\end{itemize}

\subsection{RQ1: \projName's Ability To Unlock Solver-Hostile Paths}

\begin{table}[t!]
\centering
\caption{Results on the FDLibM benchmark.
Left: KLEE baselines.
Right: LLM-based tools.
\textbf{Avg} denotes mean per-function line coverage;
\textbf{Tot} denotes overall line coverage across all functions.
Tokens are in millions.}
\label{tab:FDLibM-summary}
\footnotesize
\vspace{-2mm}
\begin{tabular}{@{}p{0.28\columnwidth}
@{\hspace{0.00\columnwidth}}
p{0.70\columnwidth}@{}}

% ================= LEFT: KLEE =================
\raggedright
\centering\textbf{KLEE Baselines}\par\vspace{2pt}
\centering
\begin{tabular}{|l|cc|}
\toprule
\multicolumn{3}{|c|}{\rule{0pt}{3.2ex}} \\   % empty row WITH vertical borders
\textbf{Solver}
& \textbf{Avg}
& \textbf{Tot} \\
\midrule
STP        & 62.68 & 70.21 \\
Z3         & 63.09 & 68.98 \\
Float-Z3   & 56.67 & 57.32 \\
\bottomrule
\end{tabular}

&
% ================= RIGHT: LLM TOOLS =================
\centering\textbf{LLM-Based Tools}\par\vspace{2pt}
\centering
\begin{tabular}{|l|ccc|ccc|}
\toprule
\multicolumn{1}{|c|}{\textbf{Backend}}
& \multicolumn{3}{c|}{\textbf{\projName{}}}
& \multicolumn{3}{c|}{\textbf{ConcoLLMic}} \\
\cmidrule(lr){2-4}\cmidrule(lr){5-7}
& \textbf{Avg}
& \textbf{Tot}
& \makecell{\textbf{Tokens}}
& \textbf{Avg}
& \textbf{Tot}
& \makecell{\textbf{Tokens}} \\
\midrule
Claude 3.7
& 91.14 & 83.62 & 0.49
& 40.02 & 31.14 & 16.45 \\

GPT-5
& \textbf{91.67} & \textbf{90.23 } & 0.54
& 38.65 & 30.62 & 10.53 \\

DeepSeek-Chat
& 91.12 & 83.31 & 0.50
& 37.17 & 28.98 & 8.78 \\
\bottomrule
\end{tabular}
\end{tabular}
\vspace{-4mm}
\end{table}


We compared \projName{} against KLEE baselines and ConcoLLMic on LogicBombs, which is designed to exercise diverse symbolic execution challenges. Each of its 53 programs has two critical paths---a logic-bomb path and a benign one---and we report bombs triggered and critical paths covered.

Table~\ref{tab:logic-bombs-summary} reports the results. KLEE triggers only 32.1--39.6\% of bombs and covers at most 56.6\% of critical paths, since solver-hostile constructs (non-linear arithmetic, floating-point conditions, complex data dependencies) leave many bomb paths unexplored. ConcoLLMic raises these to 62.3--66.0\% and 76.4--78.3\% by reasoning over execution semantics rather than strict constraints. \projName{} performs best, triggering up to 94.3\% of bombs and covering 97.2\% of critical paths---a 138.1\%/71.7\% gain over the best KLEE variant and 42.9\%/24.1\% over the best ConcoLLMic run.

%% This gain stems from \projName{}'s hybrid design: instead of replacing the SMT solver
%% or delegating entire path reasoning to an LLM, \projName{} confines LLM assistance to
%% localized, solver-hostile fragments and integrates the resulting ghost code back
%% into a standard symbolic execution pipeline.
%% As a result, \projName{} preserves global solver reasoning while selectively bypassing
%% the precise bottlenecks that block exploration of bomb-triggering paths.



% \begin{table}[h!]
% \centering
% \begin{tabular}{lccc}
% \toprule
% \textbf{Metric} &
% \textbf{KLEE (STP)} &
% \textbf{KLEE (Z3)} &
% \textbf{KLEE-FLOAT (Z3)} \\
% \midrule
% Total Items & 53 & 53 & 53 \\
% Triggered Both & 20 & 17 & 21 \\
% Triggered Only No Bomb & 19 & 22 & 18 \\
% No Trigger  & 14 & 14 & 14 \\
% Covered Critical Paths & 59 & 56 & 60 \\
% Total Critical Paths & 106 & 106 & 106 \\
% Percentage (\%) & 53.77 & 50.94 & 54.72 \\
% \bottomrule
% \end{tabular}
% \caption{Critical Path Coverage for KLEE Configurations}
% \end{table}


% \section{Results for \projName{} on the Logic Bombs Dataset}
% \label{sec:\projName{}-logic-bombs}

% Table~\ref{tab:\projName{}_logic_bombs} reports the preliminary evaluation results of \textsc{\projName{}} on the Logic Bombs dataset using three different LLM backends: Claude-3.7-Sonnet, GPT-5, and DeepSeek V3.2. The dataset contains 53 programs with a total of 106 critical paths. A program is considered to contribute two covered critical paths if both the bomb and non-bomb paths are triggered, and one covered critical path if only the non-bomb path is triggered. Formally, the number of covered critical paths is computed as $2 \times \textit{Triggered Both} + \textit{Triggered Only No Bomb}$.

% \begin{table}[ht!]
% \centering
% \caption{Results for \textsc{\projName{}} on the Logic Bombs Dataset}
% \label{tab:\projName{}_logic_bombs}
% \begin{tabular}{lccc}
% \toprule
% \textbf{Metric} 
% & \textbf{Claude-3.7-Sonnet} 
% & \textbf{GPT-5} 
% & \textbf{DeepSeek V3.2} \\
% \midrule
% Total Items & 53 & 53 & 53 \\
% Triggered Both & 50 & 49 & 45 \\
% Triggered Only No Bomb & 5 & 6 & 10 \\
% No Trigger & 0 & 0 & 0 \\
% \midrule
% Covered Critical Paths & 101 & 100 & 96 \\
% Total Critical Paths & 106 & 106 & 106 \\
% Coverage (\%) & 95.28 & 94.34 & 90.57 \\
% \bottomrule
% \end{tabular}
% \end{table}

% \subsection{Solution Strategy Breakdown for the Best-Performing Model}
% \label{sec:\projName{}-strategy-breakdown}

% For the best-performing backend, \textbf{Claude-3.7-Sonnet}, we further analyze how the 48 programs in the \textit{Triggered Both} category were solved across different reasoning strategies. The breakdown is as follows: 17 programs were solved using pure \textit{KLEE-based execution}, 12 using the \textit{model-driven abstraction}, 10 using the \textit{partition-based reasoning} strategy, and 9 using the \textit{inverse reasoning} strategy. This distribution highlights that while symbolic execution alone remains a strong baseline, a substantial portion of the hardest logic bombs require higher-level semantic reasoning enabled by LLM-driven strategies.

% \begin{table}[t]
% \centering
% \caption{Solution Strategy Distribution for Claude-3.7-Sonnet (Triggered Both = 48)}
% \label{tab:\projName{}_strategy_breakdown}
% \begin{tabular}{lc}
% \toprule
% \textbf{Strategy} & \textbf{Number of Programs} \\
% \midrule
% KLEE-only Execution & 19 \\
% Model-driven Abstraction & 12 \\
% Partition-based Reasoning & 10 \\
% Inverse Reasoning & 9 \\
% \midrule
% Total (Triggered Both) & 48 \\
% \bottomrule
% % \end{tabular}
% % \end{table}

% \section{Merged Results for all methods for Logic Bombs benchmark}
% \begin{table}[ht!]
% \centering
% \caption{Unified Critical Path Coverage Results Across Tools}
% \label{tab:merged_logic_bombs}
% \small
% \begin{tabular}{lccccccc}
% \toprule
% \textbf{Method}
% & \makecell{\textbf{Total}\\\textbf{Items}}
% & \makecell{\textbf{Trig.}\\\textbf{Both}}
% & \makecell{\textbf{Only}\\\textbf{No Bomb}}
% & \makecell{\textbf{No}\\\textbf{Trigger}}
% & \makecell{\textbf{Cov.}\\\textbf{CP}}
% & \makecell{\textbf{Total}\\\textbf{CP}}
% & \makecell{\textbf{Cov.}\\\textbf{(\%)}} \\
% \midrule
% \multicolumn{8}{l}{\emph{CP = critical paths; Cov. = covered; Trig. = triggered}} \\
% \midrule
% KLEE-STP 
% & 53 & 20 & 19 & 14 & 59 & 106 & 53.77 \\

% KLEE-Z3 
% & 53 & 17 & 22 & 14 & 56 & 106 & 50.94 \\

% KLEE-FLOAT-Z3 
% & 53 & 21 & 18 & 14 & 60 & 106 & 54.72 \\

% \midrule
% ConcoLLMic-DeepSeek 
% & 53 & 34 & 13 & 5 & 92 & 106 & 86.79 \\

% ConcoLLMic-GPT-5 
% & 53 & 33 & 15 & 5 & 91 & 106 & 85.85 \\

% ConcoLLMic-Claude-3.7 
% & 53 & 34 & 14 & 5 & 92 & 106 & 86.79 \\

% \midrule
% \projName{}-DeepSeek 
% & 53 & 43 & 10 & 0 & 96 & 106 & 90.57 \\

% \projName{}-GPT-5 
% & 53 & 47 & 6 & 0 & 100 & 106 & 94.34 \\

% \projName{}-Claude-3.7 
% & 53 & 48 & 5 & 0 & 101 & 106 & 95.28 \\
% \bottomrule
% \end{tabular}
% \end{table}

\begin{wrapfigure}{l}{0.32\textwidth}
    \vspace{-4mm}
    \centering
    \includegraphics[width=\linewidth]{venns_3.png}
    \vspace{-6mm}
    \caption{\rev{Overlap in bombs KLEE, \projName and ConcoLLMic trigger.\label{fig:bomb-venn}}}
    \vspace{-35mm}
\end{wrapfigure}

To trace this gain, Figure~\ref{fig:technique-queries} breaks coverage down by component, using Claude for the three ghost code techniques (\emph{Modeling}, \emph{Inversion}, \emph{Partition}). Each triggers a substantial but distinct subset of bombs, and their union (green row, equivalent to \projName{}) covers nearly all of them; their intersection is small, so few bombs need all three. The KLEE variants, in contrast, largely overlap. Adding KLEE to the ghost code union yields no further coverage---\projName{} is strictly additive over KLEE.

\rev{Figure~\ref{fig:bomb-venn} shows the overlap in bombs each approach triggers. The three KLEE solver variants agree closely, sharing 15 bombs with only marginal solver-specific differences at the margins. Both LLM-based approaches substantially extend KLEE's reach: all share common 18 bombs, beyond which \projName{} and ConcoLLMic jointly solve 17 more, and \projName{} exclusively solves 9. 
  The 50 bombs solved by \projName{} form a strict superset of those solved by KLEE and ConcoLLMic combined, showing both that LLM-augmented techniques extend traditional symbolic execution and that \projName{} covers all bombs reached by the other evaluated approaches on this benchmark.}

\begin{tcolorbox}[
    colback=gray!5,
    colframe=black!20,
    arc=2pt,
    boxrule=0.3mm,
    sharp corners,
    left=2pt,
    right=2pt,
    top=2pt,
    bottom=2pt,
]
\textbf{RQ1:} \projName{}'s LLM-generated ghost code unlocks solver-hostile paths that are unreachable by prior approaches, triggering up to \textbf{94.3\%} of logic bombs and covering \textbf{97.2\%} of critical paths, strictly subsuming baseline-triggered bombs. This corresponds to a \textbf{42.9\%} improvement over the strongest LLM-based baseline and a \textbf{138.1\%} improvement over the best KLEE configuration.
\end{tcolorbox}


% \wrapfill

%% \paragraph{Model sensitivity.}
%% Across all evaluated methods, differences between Claude~3.7, GPT-5, and
%% DeepSeek-Chat are small.
%% This effect is particularly pronounced for ConcoLLMic: all three models trigger
%% nearly the same number of logic bombs (34, 33 and 34 respectively), indicating that ConcoLLMic has reached a saturation point on this
%% benchmark.
%% The slightly lower result for GPT-5 is attributable to its slower per-query
%% latency under a fixed time budget rather than weaker reasoning ability.
%% Indeed, the near-identical outcomes suggest that increasing the time budget would
%% not substantially improve coverage, as the remaining bombs require structured
%% reasoning over complex paths that current LLMs cannot reliably infer when solving
%% entire path constraints.

%% Note that the Logic Bombs benchmark is particularly favorable to ConcoLLMic, because the programs are small and can be provided to the LLM in full, which reduces context fragmentation and benefits approaches that delegate whole-path reasoning to the model. 
%% Despite this advantage, ConcoLLMic still exhibits clear saturation, reinforcing
%% that its limitations stem from intrinsic reasoning difficulty rather than
%% insufficient sampling.
%% In contrast, \projName{} confines LLM reasoning to localized fragments whose complexity
%% remains largely independent of the rest of the execution path,

\begin{table}[ht!]
\centering
\caption{Line coverage (\%) results on the structured-input programs.
Left: KLEE baselines.
Right: LLM-based tools.
Avg denotes mean per-file coverage; Tot denotes overall coverage.}
\label{tab:coverage-summary}
\footnotesize
\vspace{-4mm}
\hspace*{-8mm}\begin{tabular}{@{}p{0.36\columnwidth}
@{\hspace{0.00\columnwidth}}
p{0.56\columnwidth}@{}}

% ================= LEFT: KLEE =================
\raggedright
\centering\textbf{KLEE Baselines}\par\vspace{2pt}
\centering
\begin{tabular}{|l|lcc|}
\toprule
\multicolumn{4}{|c|}{\rule{0pt}{3.2ex}} \\  % empty header row
\textbf{Bench} & \textbf{Solver} & \textbf{Avg} & \textbf{Tot} \\
\midrule
\multirow{3}{*}{libexpat}
 & Z3        & 20.21 & 14.35 \\
 & STP       & 18.10 & 12.62 \\
 & Float-Z3  & 18.40 & 12.89 \\
\midrule
\multirow{3}{*}{jq}
 & Z3        & 34.29 & 28.95 \\
 & STP       & 35.09 & 29.80 \\
 & Float-Z3  & 30.81 & 25.34 \\
\midrule
\multirow{3}{*}{bc}
 & Z3        & 50.69 & 40.01 \\
 & STP       & 54.59 & 44.70 \\
 & Float-Z3  & 18.01 & 9.84 \\
  
 \midrule
% add colour

\multirow{3}{*}{\rev{libyaml}}
 & \rev{Z3}        & \rev{21.03} & \rev{21.17} \\
 & \rev{STP}       & \rev{23.67} & \rev{23.74} \\
 & \rev{Float-Z3}  & \rev{15.18} & \rev{13.85} \\
\midrule
\midrule
\multirow{3}{*}{\rev{Overall}}
& \rev{Z3}        & \rev{32.99} & \rev{24.97} \\
& \rev{STP}       & \rev{34.31} & \rev{26.05} \\
& \rev{Float-Z3}  & \rev{20.38} & \rev{13.37} \\
% \multirow{3}{*}{libyaml}
%  & Z3        & 21.03 & 21.17 \\
%  & STP       & 23.67 & 23.74 \\
%  & Float-Z3  & 15.18 & 13.85 \\

% \midrule
% \multirow{3}{*}{Overall}
% & Z3        & 32.99 & 24.97 \\
%  & STP       & 34.31 & 26.05 \\
%  & Float-Z3  & 20.38 & 13.37 \\
 % old
 % & Z3        & 36.97 & 25.43 \\
 % & STP       & 37.86 & 26.33 \\
 % & Float-Z3  & 22.11 & 13.31 \\
\bottomrule
\end{tabular}

&
% ================= RIGHT: LLM TOOLS =================
\centering\textbf{LLM-Based Tools}\par\vspace{2pt}
\centering
\begin{tabular}{|l|cc|cc|cc|}
\toprule
\multicolumn{1}{|c|}{\textbf{Backend}}
 & \multicolumn{2}{c|}{\textbf{\projName{}}}
 & \multicolumn{2}{c|}{\textbf{ConcoLLMic}}
 & \multicolumn{2}{c|}{\textbf{Cottontail}} \\
\cmidrule(lr){2-3}\cmidrule(lr){4-5}\cmidrule(lr){6-7}
 & \textbf{Avg} & \textbf{Tot}
 & \textbf{Avg} & \textbf{Tot}
 & \textbf{Avg} & \textbf{Tot} \\
\midrule
Claude 3.7    & 46.29 & 38.63 & 11.27 & 4.18 & 11.42 & 11.38 \\
GPT-5         & 59.72 & \textbf{47.75} & 13.04 & 5.11 & 21.86 & 18.30 \\
DeepSeek-Chat & 49.68 & 43.03 & 8.80 & 2.37 & 12.32 & 13.13 \\
\midrule
Claude 3.7    & 46.19 & \textbf{45.37} & 9.68 & 8.37 & 30.56 & 29.37 \\
GPT-5         & 48.69 & 40.98 & 9.75 & 8.56 & 35.34 & 31.74 \\
DeepSeek-Chat & 56.49 & 44.86 & 12.69 & 10.00 & 31.54 & 30.38 \\
\midrule
Claude 3.7    & 70.74 & 59.89 & 15.23 & 6.37 & 52.10 & 39.46 \\
GPT-5         & 79.15 & \textbf{69.79} & 17.13 & 6.87 & 56.36 & 48.12 \\
DeepSeek-Chat & 65.65 & 62.93 & 11.83 & 4.65 & 54.03 & 46.09 \\

\midrule
\rev{Claude 3.7}    & \rev{58.03} & \rev{56.51} & \rev{29.74} & \rev{29.17} & \rev{53.5} & \rev{52.43} \\
\rev{GPT-5}         & \rev{63.40} & \rev{\textbf{65.59}} & \rev{27.53} & \rev{26.33} & \rev{44.75} & \rev{43.86} \\
\rev{DeepSeek-Chat} & \rev{62.36} & \rev{61.72} & \rev{29.10} & \rev{26.96} & \rev{55.35} & \rev{55.24} \\
\midrule
\midrule
\rev{Claude 3.7}    & \rev{56.75} & \rev{48.20} & \rev{16.76} & \rev{8.05} & \rev{38.73} & \rev{26.91} \\
\rev{GPT-5}         & \rev{64.21} & \rev{\textbf{56.07}} & \rev{17.20} & \rev{8.35} & \rev{41.21} & \rev{32.20} \\
\rev{DeepSeek-Chat} & \rev{59.29} & \rev{51.70} & \rev{15.67} & \rev{6.61} & \rev{40.15} & \rev{30.72} \\
% \midrule
% Claude 3.7    & 58.03 & 56.51 & 29.74 & 29.17 & 53.5 & 52.43 \\
% GPT-5         & 63.40 & \textbf{65.59} & 27.53 & 26.33 & 44.75 & 43.86 \\
% DeepSeek-Chat & 62.36 & 61.72 & 29.10 & 26.96 & 55.35 & 55.24 \\
% \midrule
% Claude 3.7    & 56.75 & 48.20 & 16.76 & 8.05 & 38.73 & 26.91 \\
% GPT-5         & 64.21 & \textbf{56.07} & 17.20 & 8.35 & 41.21 & 32.20 \\
% DeepSeek-Chat & 59.29 & 51.70 & 15.67 & 6.61 & 40.15 & 30.72 \\
% old  
% Claude 3.7    & 56.33 & 47.19 & 12.43 & 5.49 & 33.80 & 23.82 \\
% GPT-5         & 64.48 & \textbf{54.91} & 13.76 & 6.17 & 40.03 & 30.79 \\
% DeepSeek-Chat & 58.26 & 50.49 & 11.19 & 4.14 & 35.15 & 27.25 \\
\bottomrule
\end{tabular}

\end{tabular}
\vspace{-4mm}
\end{table}


\subsection{RQ2: \projName's Applicability To Real-World Projects}

We evaluate \projName{} on FDLibM and the structured-input programs separately, since they involve different baselines.
Table~\ref{tab:FDLibM-summary} reports line coverage on FDLibM, where coverage reflects each technique's ability to explore floating-point-guarded control flow rather than reach explicit bomb branches. KLEE achieves 56.7--63.1\% average and up to 70.2\% total coverage; STP and Z3 perform comparably, with STP marginally ahead due to higher solver throughput. KLEE-Float-Z3 trails despite more precise floating-point reasoning, for two reasons: slow float solving limits exploration (only 21.37\% as many test cases as standard KLEE), and its 9-year-old KLEE fork lacks modern performance optimizations.

ConcoLLMic, in contrast to its LogicBombs results, performs poorly on FDLibM (37.2--40.0\% average, 29.0--31.1\% total coverage). Even its best run with Claude-3.7 yields 55.66\% less coverage than KLEE-STP despite issuing up to 16.45M tokens. This exposes a fundamental weakness of whole-path LLM reasoning on numerical code: predicting inputs that satisfy long chains of floating-point constraints is unreliable from semantic reasoning alone.

\projName{} achieves the highest coverage on FDLibM, up to 91.7\% average and 90.2\% total---an average 183.5\% improvement over ConcoLLMic and 22.09\% over the best KLEE version, despite generating only 13.87\% as many test cases. To rule out better time-budget allocation as the explanation, we combine the coverage of all three KLEE variants and all three ConcoLLMic configurations, a baseline with 6$\times$ the exploration time of a single \projName{} run. This union reaches only 70.75\% average and 77.95\% total coverage, while each individual \projName{} configuration exceeds it by 6.88--15.75\% in total coverage. \projName{}'s gains thus come from reaching solver-hostile branches that remain inaccessible even when the strengths of existing approaches are aggregated.

% \midrule
% \multirow{3}{*}{Avg.}
%  & Z3        & 35.06 & 27.77 \\
%  & STP       & 35.93 & 29.04 \\
%  & Float-Z3  & 22.42 & 16.02 \\
% \midrule
% Claude 3.7   & 54.41 & 47.96 & 12.06 & 6.31 & 31.36 & 26.74 \\
% GPT-5        & 62.52 & \textbf{52.84} & 13.31 & 6.85 & 37.85 & 32.72 \\
% DeepSeek-Chat & 57.27 & 50.27 & 11.11 & 5.67 & 32.63 & 29.87 \\



% \begin{table}[ht!]
% \centering
% \small
% \caption{Total token usage in millions of tokens across structured-input benchmarks.}
% \label{tab:token-usage}

% \begin{tabular}{|l|l|ccc|}
% \toprule
% % \multicolumn{5}{|c|}{\rule{0pt}{3.2ex}} \\  % empty header row for visual alignment
% \textbf{Bench} & \textbf{Model}
%  & \textbf{\projName{}}
%  & \textbf{ConcoLLMic}
%  & \textbf{Cottontail} \\
% \midrule
% \multirow{3}{*}{libexpat}
%  & Claude 3.7    & 0.15 & 5.78 & 1.15 \\
%  & GPT-5         & 0.17 & 4.01 & 0.71 \\
%  & DeepSeek-Chat & 0.15 & 3.75 & 1.19 \\
% \midrule
% \multirow{3}{*}{jq}
%  & Claude 3.7    & 0.07 & 2.79 & 2.08 \\
%  & GPT-5         & 0.08 & 1.74 & 1.02 \\
%  & DeepSeek-Chat & 0.06 & 1.57 & 1.53 \\
% \midrule
% \multirow{3}{*}{bc}
%  & Claude 3.7    & 0.14 & 10.34 & 2.34 \\
%  & GPT-5         & 0.17 & 8.30 & 0.85 \\
%  & DeepSeek-Chat & 0.15 & 7.93 & 0.99 \\
% \midrule
% \multirow{3}{*}{Overall}
%  & Claude 3.7    & 0.36 & 18.91 & 5.57 \\
%  & GPT-5         & 0.42 & 14.05 & 2.58 \\
%  & DeepSeek-Chat & 0.36 & 13.25 & 3.71 \\
% \bottomrule
% \end{tabular}
% \end{table}

% \multirow{3}{*}{Average.}
%  & Claude 3.7    & 0.23 & 7.88 & 16.37 \\
%  & GPT-5         & 0.26 & 5.34 & 0.86 \\
%  & DeepSeek-Chat & 0.24 & 5.01 & 1.24 \\
%  \midrule

Table~\ref{tab:coverage-summary} reports line coverage on the structured-input benchmarks, which exercise grammar-constrained parsing and semantic validation. KLEE variants achieve modest coverage, with total coverage below 40\% even for the best configuration. Z3 and STP perform comparably, while Float-Z3 trails sharply---most notably on \texttt{bc}, where missing inline-assembly support limits it to a single test case---and elsewhere generates under 1\% as many tests as standard KLEE.

Neither ConcoLLMic nor Cottontail consistently outperforms KLEE on structured-input programs. ConcoLLMic falls below 10\% total coverage---only \rev{29.44\%} of KLEE-STP's coverage on average---largely due to instrumentation failures: functions exceeding LLM output limits, and instrumentation points that cannot safely accommodate injected logging, both yielding zero coverage for the affected files. Cottontail fares better since its grammar-based generation matches structured-input programs naturally, beating the best KLEE variant by \rev{23.61\%} with GPT-5; it is strongest on \texttt{libyaml} and \texttt{bc} but struggles on \texttt{libexpat}, where coverage requires deep semantic validation beyond syntactic correctness. Overall, end-to-end LLM-driven concolic execution remains fragile in the presence of complex cross-function dependencies.

\projName{} outperforms all baselines on every structured-input benchmark and every backend, improving total coverage by \rev{85.03--115.24\%} over the strongest KLEE, by \rev{6.8$\times$} on average over ConcoLLMic, and by \rev{83.36\% on average over Cottontail (29.94\% $\to$ 51.9\% total coverage)}.

\begin{tcolorbox}[colback=gray!5, colframe=black!20, arc=2pt, boxrule=0.3mm, breakable, sharp corners, left=2pt, right=2pt, top=2pt, bottom=2pt]
\textbf{RQ2:} \projName{} scales  to real-world numerical and structured-input programs, consistently outperforming symbolic execution and LLM-based baselines.
On FDLibM, \projName{} improves coverage by \textbf{183.5\%} on average over ConcoLLMic and by \textbf{22.09\%} over the strongest KLEE configuration.
On structured-input programs, \projName{} improves total coverage by up to \textbf{\rev{115.24\%}} over KLEE, by \textbf{\rev{83.36\%}} on average over Cottontail, and covers \textbf{\rev{6.8}$\times$} more lines than ConcoLLMic.
\end{tcolorbox}


\begin{revblock}
%% \subsection{Analysis of Generated Inversions}

\begin{table}[t]
\centering
\begin{minipage}[b]{0.48\linewidth}
\centering
% \begin{tabular}{lr}
% \hline
% \rev{Stage} & \rev{Percentage (\%)} \\
% \hline
% \rev{Compiled successfully} & \rev{93} \\
% \rev{Produced KTests} & \rev{85} \\
% \rev{Usable inversion} & \rev{63} \\
% \rev{New coverage achieved} & \rev{45} \\
% \hline
% \end{tabular}

{\footnotesize
\begin{tabular}{lr}
\hline
\rev{Outcome category} & \rev{Attempts (\%)} \\
\hline
\rev{A1-Generation/execution failure} & \rev{7} \\
\rev{A2-Interface mismatch } & \rev{8} \\
\rev{A3-Semantically invalid inversion} & \rev{22} \\
\rev{P1-Non convergence} & \rev{18} \\
\rev{P2-Successful path trigger} & \rev{45} \\
\hline
\end{tabular}
}
\caption{\rev{Mutually exclusive categories of the outcomes for 100 inversion attempts from LogicBombs.}}
\label{tab:inversion_pipeline}
\end{minipage}
\hfill
\begin{minipage}[b]{0.48\linewidth}
\centering
{\footnotesize
\begin{tabular}{lcc}
\hline
\rev{Type} & \rev{\makecell{Additional \\ Coverage (\%)}} & \rev{\makecell{Avg. \\ Attempts}} \\
\hline
\rev{Partition} & \rev{93.2} & \rev{1.2} \\
\rev{Modeling} & \rev{62.1} & \rev{1.4} \\
\rev{Inversion} & \rev{43.2} & \rev{2.3} \\
\hline
\end{tabular}
}
\caption{\rev{Effectiveness of ghost-code strategies across LogicBombs, fdlibm, structured inputs benchmarks.}}
\label{tab:strategy_effectiveness}
\end{minipage}
\label{tab:inversion_results}
\vspace{-8mm}
\end{table}

We analyze 100 randomly sampled inversion attempts from LogicBombs, excluding malformed outputs such as truncated or unparsable responses, classifying their first failure stage (Table~\ref{tab:inversion_pipeline}). The first three categories are \emph{artifact defects}; the last two are at least approximations of the true inverse.

\begin{description}\itemsep1pt
\item[(A1) Generation/execution failure (7\%).] Does not compile, crashes, or cannot run.
\item[(A2) Interface mismatch (8\%).] Runs but output is unusable by Gordian's refinement (missing variables, wrong type, schema violation).
\item[(A3) Semantically invalid (22\%).] Correct interface, but not a valid inverse or approximation.
\item[(P1) Non-convergence (18\%).] Refinement does not converge within budget. Manual inspection suggests the candidates are plausible approximations; the loop fails to close mainly due to infeasible target paths or approximation error too large for the optimization step to absorb, not inversion defects. This is the operational \textsc{Unsat} case of \Cref{sec:bidirectional}.
\item[(P2) Successful path trigger (45\%).] The inversion yields a concrete input validated on the original program that reaches the uncovered path.
\end{description}

Two rates summarize the distribution. The artifact-level success rate (defect-free inversions, A1--A3 excluded) is $63\%$, with $93\%$ compiling and $85\%$ executing. The end-to-end rate is $45\%$; the $18\%$ gap is category P1, where the inversion is plausible but the refinement loop does not close---a pipeline property, not an inversion defect. Many successful inversions are themselves approximate (search-based or heuristic) rather than exact, yet propagate values consistent enough for refinement to converge. LLMs are thus reliable enough as inversion generators to aid symbolic execution.


Across FDLibM and the structured-input benchmarks, \projName{} identifies one solver-hostile fragment per \(117.5\) executable lines on average, targeted with one or more of the three ghost-code strategies (Table~\ref{tab:strategy_effectiveness}). Partitioning is the most effective, achieving the highest coverage with minimal attempts. Modeling balances generality and effectiveness. Inversion is less consistent and needs more attempts, but is crucial in solver-hostile numerical settings where other strategies fail.
\end{revblock}

\subsection{RQ3: \projName's Token Efficiency}

% \begin{wraptable}{r}{0.55\textwidth}
% \centering
% \small
% \caption{Total token usage (in Millions) across structured-input benchmarks.}
% \label{tab:token-usage}
% \begin{tabular}{lccc}
% \toprule
% Model & \projName{} & ConcoLLMic & Cottontail \\
% \midrule
% Claude 3.7    & 0.36 & 18.91 & 5.57 \\
% GPT-5         & 0.42 & 14.05 & 2.58 \\
% DeepSeek-Chat & 0.36 & 13.25 & 3.71 \\
% \bottomrule
% \end{tabular}
% \end{wraptable}

\projName{} uses the LLM fundamentally differently from ConcoLLMic and Cottontail: it invokes the LLM only at specific solver-hostile fragments to generate a small amount of ghost code, then delegates the rest to the symbolic executor (\rev{LLM calls account for only ${\sim}9.2\%$ of Gordian's execution time}). ConcoLLMic and Cottontail, by contrast, query the LLM continuously to reason about execution paths and synthesize inputs, incurring LLM cost throughout exploration.

\begin{wraptable}{r}{0.5\textwidth}
  \vspace{-3mm}
\centering
\caption{Total token usage (in Millions) across structured-input benchmarks.}
\label{tab:token-usage}
\footnotesize
\begin{tabular}{lccc}
\toprule
\rev{Model} & \rev{\projName{}} & \rev{ConcoLLMic} & \rev{Cottontail} \\
\midrule
\rev{Claude 3.7}    & \rev{0.40} & \rev{24.41} & \rev{6.87} \\
\rev{GPT-5}         & \rev{0.47} & \rev{14.90} & \rev{2.81} \\
\rev{DeepSeek-Chat} & \rev{0.40} & \rev{14.91} & \rev{5.02} \\
\bottomrule
\end{tabular}
  \vspace{-2mm}
\end{wraptable}

On LogicBombs (\Cref{tab:logic-bombs-summary}), \projName{} uses roughly $10\times$ fewer tokens than ConcoLLMic (0.32M vs.\ 3.07M per run on average). On FDLibM (\Cref{tab:FDLibM-summary}), the reduction grows to $23\times$. On the structured-input benchmarks (Table~\ref{tab:token-usage}), ConcoLLMic consumes on average \rev{42.69}$\times$ more tokens than \projName{}, while Cottontail consumes between \rev{5.98}$\times$ (GPT-5) and \rev{17.17}$\times$ (Claude-3.7). \projName{} thus achieves higher coverage at substantially lower LLM cost than existing LLM-based testing approaches.

% gordian 3.76
% concollmic 99.19
\begin{tcolorbox}[colback=gray!5, colframe=black!20, arc=2pt, boxrule=0.3mm,
                  breakable, sharp corners, left=2pt, right=2pt, top=2pt, bottom=2pt]
\textbf{RQ3:} \projName{} achieves more than an order-of-magnitude reduction
in LLM token cost while achieving higher coverage.
Across all 5 benchmarks, ConcoLLMic consumes \textbf{\rev{26.38}$\times$} as many tokens
as \projName{} (a \textbf{\rev{96.21}\%} reduction), while
Cottontail consumes \textbf{\rev{11.57}$\times$} as many tokens
(a \textbf{\rev{91.36}\%} reduction for \projName{}).
\end{tcolorbox}



\begin{revblock}
\subsection{RQ4: \projName{}'s Coverage Progress Over Time}

To assess whether \projName{}'s gains are realized only after long runs, we measure coverage as a function of wall-clock time across all evaluated tools.

\begin{figure}[t]
    \centering
    \includegraphics[width=0.49\textwidth]{curves_structured.png}
    \hfill
    \includegraphics[width=0.49\textwidth]{curves_fdlibm.png}
    \vspace{-7mm}
    \caption{\begin{revblock}
        Coverage progression over time for structured-input programs (left) and FDLibM (right).
    \end{revblock}}
    \label{fig:coverage-curves}
    \vspace{-4mm}
\end{figure}

Figure~\ref{fig:coverage-curves} shows coverage versus time budget. All approaches follow a common pattern---rapid initial gains followed by a plateau---but differ in where their gains concentrate. KLEE begins exploring immediately and accrues coverage early, while LLM-based approaches incur an initial delay for preprocessing and LLM calls; once exploration starts, their growth trends resemble KLEE's.

\projName{}'s curves are distinctly step-wise: large early jumps correspond to successful ghost-code transformations (partitioning, modeling) that unlock previously unreachable regions, and subsequent improvements occur in discrete steps as new partitions or solver-aiding transformations are introduced. \projName{} typically surpasses both KLEE and ConcoLLMic well within the first few hours, and the bulk of its advantage over KLEE is realized before the long-budget plateau---meaning the gains are not an artifact of the 36-hour ceiling. Cottontail and ConcoLLMic, by contrast, grow more smoothly because their LLM interactions are interleaved continuously with exploration.

On FDLibM, KLEE quickly generates floating-point-heavy test cases for high initial coverage, but stagnates early as many paths hit solver-hostile constraints it cannot extend. \projName{} continues to make progress by resolving such bottlenecks via targeted transformations. Minor spikes in the KLEE curves are a measurement artifact: coverage updates only when test cases are generated, and several paths sharing prefixes can reach new code within a short window.

\begin{tcolorbox}[colback=gray!5, colframe=black!20, arc=2pt, boxrule=0.3mm, breakable, sharp corners, left=2pt, right=2pt, top=2pt, bottom=2pt]
\rev{\textbf{RQ4:} \projName{}'s advantage emerges early rather than only at long budgets. Many of its gains over KLEE and ConcoLLMic appear within the first 20\% of the 36-hour budget, while additional gains continue later through step-wise improvements unlocked by ghost-code transformations.}
\end{tcolorbox}
\end{revblock}


\begin{revblock}
    
\subsection{RQ5: \projName{}'s Exploration of Hard-to-Reach Paths}
\label{sec:hard-paths}

To check whether \projName{}'s gains correspond to genuinely new paths rather than redundant coverage, we analyze the overlap of covered lines across techniques. Coverage is a standard testing-effectiveness proxy in symbolic and concolic execution and provides more stable evaluation signals than bug-based metrics~\cite{FSE26-concordance}; the overlap view further reveals how much coverage is shared and how much each technique contributes uniquely.

\begin{figure}[t]
\centering
\includegraphics[width=0.9\textwidth]{upset_2.png}

\caption{\begin{revblock}
Overlap of covered lines across techniques on \texttt{libyaml}, aggregated across all models.
The largest intersection corresponds to lines covered exclusively by \projName{} (19.2\% of all covered lines), exceeding even the shared \projName{}--Cottontail intersection (17.8\%).
\end{revblock}}
\label{fig:libyaml-upset}
\end{figure}

Figure~\ref{fig:libyaml-upset} shows the overlap on \texttt{libyaml}, aggregated across LLM backends (the best-performing structured-input benchmark for LLM-based approaches). Three observations stand out: \projName{} achieves the largest overall coverage; it contributes a substantial set of lines covered by no other approach, reaching previously unexplored behaviors; and although it overlaps significantly with structure-aware Cottontail, it still extends coverage beyond their union---its gains are not reducible to input-structure modeling alone.

\begin{tcolorbox}[colback=gray!5, colframe=black!20, arc=2pt, boxrule=0.3mm, breakable, sharp corners, left=2pt, right=2pt, top=2pt, bottom=2pt]
\rev{\textbf{RQ5:} \projName{} reaches hard coverage regions not attained by the other evaluated techniques. On \texttt{libyaml}, \textbf{19.2\%} of covered lines are reached by \projName{} alone---exceeding even its overlap with Cottontail---so its advantage cannot be explained by input-structure modeling or redundant coverage of paths reachable by other techniques.}
\end{tcolorbox}
\end{revblock}
    

\begin{revblock}
  \subsection{Rerun Variability Analysis.}
  
To quantify the impact of LLM stochasticity, we evaluate a representative subset of 16 source files (around 10\% of the total), uniformly sampled across benchmarks. For each file and each of the three LLM backends, we perform three independent reruns, yielding 48 benchmark–model pairs and 144 executions in total. 

We measure variability using the within-pair standard deviation across reruns. The average within-pair standard deviation is 3.08 coverage points overall, with per-model values of 2.65 (Claude), 2.36 (GPT), and 4.22 (DeepSeek), indicating moderate stochasticity at the level of individual runs.

To assess the impact on results, we approximate the standard error of average coverage over $N$ files as $\mathrm{SE} \approx s / \sqrt{N}$. This yields an approximate 95\% stochasticity band of ±0.68 coverage points for FDLibM (78 files; ±0.59 for Claude, ±0.52 for GPT, and ±0.94 for DeepSeek) and ±1.38 points for structured-input programs (19 files; ±1.19 for Claude, ±1.06 for GPT, and ±1.90 for DeepSeek). Per-model bounds remain below ±1 point for FDLibM and below ±2 points for structured benchmarks.
These deviations are small relative to the observed improvements (double digits coverage points over baselines), indicating that our findings are robust to LLM stochasticity.
\end{revblock}
    


%% \section{Motivating Example: Inverting Trigonometric Guards in Complex Logic}

%% To illustrate the strength of our approach over Concolic LLMs such as ConcoLLMic, consider a real-world logic fragment from the \texttt{FDLibM} math library implementing \texttt{the \_ieee754\_j0(x)},
%% which is Bessel function of the first kinds of order zero. The relevant logic is as follows:

%% \begin{verbatim}
%% hx = __HI(x);
%% ix = hx & 0x7fffffff;
%% if (ix >= 0x7ff00000) return one / (x * x);
%% x = fabs(x);
%% if (ix >= 0x40000000) {
%%     s = sin(x);
%%     c = cos(x);
%%     ss = s - c;
%%     cc = s + c;
%%     if (ix < 0x7fe00000) { 
%%         z = -cos(x + x);
%%         if ((s * c) < 0)
%%             cc = z / ss;
%%         else
%%             ss = z / cc;
%%     }
%% }
%% \end{verbatim}

%% This example contains a mix of low-level bitmasking logic, floating-point branching, and nonlinear trigonometric functions. For symbolic execution (e.g., via KLEE), the presence of \texttt{sin} and \texttt{cos} in guard conditions like \texttt{(s * c) < 0} causes immediate concretization, as solvers cannot handle inversion of transcendental functions. As a result, even though the preceding path (e.g., \texttt{|x| >= 2.0}) is solvable by symbolic means, the solver fails to complete the path condition.

%% To overcome this, we resymbolize the intermediate variables \texttt{s} and \texttt{c}, and target the branch:
%% \[
%% (s \cdot c) < 0
%% \]
%% We resymbolize the intermediate variables \texttt{s} and \texttt{c}, constraining their values to lie within the range $[-1, 1]$ as determined by the definitions of $\sin(x)$ and $\cos(x)$. A straightforward pair of assignments that satisfy the condition \texttt{(s * c) < 0} is \texttt{s = 1}, \texttt{c = -1}.
%%  These values are, however, not realizable simultaneously under real inputs (since $\sin^2(x) + \cos^2(x) > 1$), but we can invert them to find the closest valid input $x$.

%% We use the following Python routine to perform this inversion:

%% \begin{verbatim}
%% import math
%% from scipy.optimize import minimize_scalar

%% def invert_sin_cos_loose(s_target, c_target):
%%     def error(x):
%%         s = math.sin(x)
%%         c = math.cos(x)
%%         return (s - s_target)**2 + (c - c_target)**2
    
%%     result = minimize_scalar(error, bounds=(0, 2 * math.pi), method='bounded')
%%     return result.x if result.success else None
%% \end{verbatim}

%% When applied to \texttt{s = 1}, \texttt{c = -1}, the output is $x \approx 2.356$ radians ($135^\circ$), which best approximates the desired trigonometric behavior. Crucially, this value satisfies the prefix conditions:

%% \begin{itemize}
%%     \item \texttt{|x| >= 2.0} $\Rightarrow$ \texttt{ix >= 0x40000000}
%%     \item \texttt{x + x < overflow} $\Rightarrow$ \texttt{ix < 0x7fe00000}
%% \end{itemize}

%% Thus, the path is fully enabled. Our system delegates the difficult trigonometric inversion to the LLM (using gradient-guided estimation), while leaving range-based prefix logic and overflow guards to symbolic execution, which handles them well.

%% \paragraph{C Implementation.}
%% For integration into our execution pipeline, the LLM (gpt-5) implements the same inversion logic in C.
%% The routine reads the post-NL values of $s$ and $c$ from a KLEE-generated JSON file,
%% performs bounded numerical inversion, and emits concrete assumptions on $x$.

%% \begin{verbatim}
%% static double golden_section_minimize(double (*f)(double, void *),
%%                                       void *ctx,
%%                                       double a, double b,
%%                                       int max_iter, double tol) {
%%     const double gr = (sqrt(5.0) - 1.0) * 0.5;
%%     double c = b - gr * (b - a);
%%     double d = a + gr * (b - a);
%%     double fc = f(c, ctx), fd = f(d, ctx);

%%     for (int i = 0; i < max_iter && fabs(b - a) > tol; i++) {
%%         if (fc < fd) {
%%             b = d; d = c; fd = fc;
%%             c = b - gr * (b - a);
%%             fc = f(c, ctx);
%%         } else {
%%             a = c; c = d; fc = fd;
%%             d = a + gr * (b - a);
%%             fd = f(d, ctx);
%%         }
%%     }
%%     return 0.5 * (a + b);
%% }

%% static double sincos_error(double x, void *ctx) {
%%     double s = sin(x), c = cos(x);
%%     double ds = s - ((double *)ctx)[0];
%%     double dc = c - ((double *)ctx)[1];
%%     return ds * ds + dc * dc;
%% }

%% void build_pre_nl_assumes(const char *json_path) {
%%     double s, c;
%%     json_return_double(json_path, "s", &s);
%%     json_return_double(json_path, "c", &c);

%%     double ctx[2] = {s, c};
%%     double x = golden_section_minimize(
%%         sincos_error, ctx, 0.0, 2*M_PI, 200, 1e-12
%%     );

%%     printf("klee_assume((x - %.17g) <= 0);\n", x);
%%     printf("klee_assume((x - %.17g) >= 0);\n", x);
%% }
%% \end{verbatim}

%% \subsection*{Why ConcoLLMic Fails}

%% ConcoLLMic integrates LLM-driven reasoning into symbolic execution by delegating certain subexpressions—typically leaf-level computations—to a language model that proposes candidate inputs. It concretizes these subexpressions into queries like \texttt{f(x) == y} and attempts to guess suitable $x$ via natural language synthesis and python execution.

%% However, in this case, ConcoLLMic is limited by the structure and entanglement of the path constraints. Specifically, the branch condition \texttt{(s * c) < 0} depends on \texttt{s = sin(x)} and \texttt{c = cos(x)}, but \texttt{x} itself is not explicitly constrained at the top level. Instead, it is embedded behind several nontrivial bitwise manipulations and floating-point operations:

%% \begin{itemize}
%%     \item \texttt{hx = \_\_HI(x);} extracts the high 32 bits of the IEEE-754 representation.
%%     \item \texttt{ix = hx \& 0x7fffffff;} masks off the sign bit to obtain \texttt{|x|}.
%%     \item Guard checks like \texttt{ix >= 0x40000000} (i.e., \texttt{|x| >= 2.0}) are applied before the trigonometric terms are even computed.
%% \end{itemize}

%% ConcoLLMic would need to "guess" values for \texttt{x} that not only cause \texttt{s * c < 0} but also satisfy all the prefix and suffix constraints, including floating-point format checks and overflow prevention (\texttt{x + x < threshold}). Since the model is not a solver, and the trigonometric constraints do not appear in an isolated leaf computation but rather intertwined with program state and memory-level logic, ConcoLLMic cannot effectively predict such values.

%% In contrast, our method decouples the hard-to-invert trigonometric condition. We symbolically solve the prefix logic (bit checks, ranges), while resymbolizing \texttt{s} and \texttt{c} and inverting them with an optimizer to find the closest $x$ that satisfies the nonlinear guard. This hybrid delegation of roles—solver for prefix, LLM for inversion—enables us to trigger the path with higher reliability and scalability.

%% This example illustrates how \projName{} can trigger semantically rich paths that are unreachable to either traditional symbolic executors or purely LLM-driven concolic methods.

\section{Discussion}

Although \projName significantly outperformed baselines in the evaluation, it has several limitations. First, \projName operates by injecting LLM-generated ghost code and recompiling the target program for symbolic execution. \rev{This introduces engineering overhead, especially for large systems with complex build pipelines, external dependencies, or custom build scripts. In practice, the overhead is mitigated by incremental builds: \projName{} modifies only the harness or a small number of target/source files, so build systems such as \texttt{make} typically recompile only the affected objects and relink the executable rather than rebuilding the entire project. Nevertheless, manual integration effort may still be required for complex projects.} Second, \projName relies on general-purpose LLMs, which may reflect biases or implicit knowledge from their training data. In addition, LLM-generated ghost code may be incomplete or imprecise, particularly for inversion and semantic modeling. Importantly, \projName{} treats all LLM outputs as heuristic aids: all generated test inputs are replayed on the original program, ensuring that incorrect abstractions affect efficiency rather than soundness; however, it may still miss feasible paths. Third, \projName is most effective when solver-hostile behavior is localized to identifiable program fragments. Programs whose complexity arises from deeply entangled global invariants or hard-to-model environment interactions may benefit less from ghost code generation. Fourth, \projName's bidirectional constraint propagation algorithm in \Cref{sec:bidirectional} relies on heuristic optimization, and does not guarantee convergence.

A natural direction for future work is a tighter integration of \projName's capabilities directly inside a symbolic execution engine. In our current implementation, \projName{} operates as an external orchestrator that interacts with KLEE through repeated recompilation and execution. While effective, this architecture introduces engineering complexity and multiple points of failure. A more principled approach would embed LLM assistance natively within KLEE. Solver queries or path fragments that exceed resource limits could be detected online and selectively delegated to LLM-based reasoning modules. Once a generated model or inverse is validated by successfully triggering a previously unreachable path, it could be cached and reused across executions or program versions. Over time, this would yield an evolving library of solver-friendly models and inverses, potentially compositional and shared across analyses.

\begin{revblock}

A limitation of the bidirectional constraint propagation algorithm presented in \Cref{sec:bidirectional} is that it assumes a single difficult fragment along any explored path. Accordingly, our implementation never attempts to invert more than one fragment at a time. This is a pragmatic choice rather than a fundamental restriction: empirically, addressing even a single difficult fragment yields substantial practical gains. In principle, however, the algorithm extends naturally to paths containing multiple difficult fragments, as follows. Suppose symbolic execution of the transformed program $P'$ encounters a sequence of $n$ difficult fragments $C_1,\ldots,C_n$ along a single path, each equipped with an LLM-generated inverse $f_i^{-1}$. The accumulated path condition then decomposes into $n+1$ contiguous segments,
\[
\pi \;=\; \pi_0 \,\wedge\, \pi_1 \,\wedge\, \cdots \,\wedge\, \pi_n,
\]
where $\pi_0$ constrains only the original input variables (the prefix preceding $C_1$), and each subsequent $\pi_i$ ($1\le i\le n$) is the segment between $C_i$ and $C_{i+1}$ (with $\pi_n$ the final suffix), depending on the fresh symbolic variables used in $C_1,\ldots,C_i$ as well as on the original inputs.

The generalized procedure might proceed by ``meeting in the middle'' across all $n$ cuts simultaneously. It first solves the bottom segment $\pi_n$ to obtain an assignment $A_n$, then propagates values upward fragment by fragment: from $A_i$ it extracts the values written by $C_i$, applies $f_i^{-1}$ to obtain a suggested pre-state $\sigma_i$ over the variables read in $C_i$, and checks whether $\sigma_i$ together with $A_i$ is consistent with the segment $\pi_{i-1}$. If every $\pi_{i-1}$ is satisfied along the way and reaches $\pi_0$ consistently, a model for the whole path has been found and the algorithm terminates. Otherwise, at the highest segment $\pi_{j-1}$ that becomes unsatisfied, the algorithm invokes the optimizing solver to compute an assignment $A_{j-1}\models\pi_{j-1}$ that is as close as possible to $\sigma_j$, and then propagates $A_{j-1}$ downward by concretely executing $f_{j-1}, f_j, \ldots, f_n$ in turn. The upward--downward traversal is then repeated from the new assignment, until either a fixpoint that satisfies $\pi_0\wedge\cdots\wedge\pi_n$ is reached, or the procedure stalls and conservatively returns \textsc{Unsat}. We leave the implementation, convergence analysis, and empirical evaluation of this multi-fragment extension to future work.

\end{revblock}


%% For example, KLEE could automatically identify solver-hostile regions, trigger
%% LLM-guided semantic partitioning or inversion on demand, and resume symbolic
%% execution using the generated abstractions without recompilation.
%% Such an integration would also enable persistent reuse of verified abstractions.
%% This mechanism could substantially reduce LLM invocation cost, latency, and
%% variance, while improving robustness.
%% More broadly, this direction points toward symbolic execution engines that
%% dynamically synthesize, validate, and reuse semantic knowledge during analysis,
%% combining the rigor of constraint solving with the adaptability of learned
%% reasoning.


\section{Related Work}

\paragraph*{Symbolic and Concolic Testing} Symbolic execution~\cite{king1976symbolic} (e.g., KLEE~\cite{klee}) explores program paths over symbolic inputs, but is limited by path explosion and solver bottlenecks, especially with complex arithmetic and hard-to-model environments~\cite{env_modeling,autobug}. Concolic execution~\cite{dart} improves practicality by interleaving concrete runs with symbolic reasoning; systems such as S2E~\cite{s2e} scale this to whole-system settings, while compilation-based engines (e.g., SymCC~\cite{symcc}) improve throughput via native execution and lightweight instrumentation. Prior work addresses exploration blowups from loops and recursion (e.g., MARCO~\cite{hu2024marco}), whereas \projName targets the complementary bottleneck: solver-hostile fragments that block progress. \rev{Environment modeling via synthesis (e.g., SE-ESOC~\cite{mechtaev2018symbolic}) replaces missing or external semantics with generated code; \projName differs in that it targets internal program fragments whose code is present but solver-hostile, and uses ghost code not merely as a stub but as a solver aid for propagating and reconciling constraints across the original fragment. \projName generalizes this idea by using LLMs to generate ghost code that (i) propagates constraints through hard fragments, (ii) introduces solver-friendly surrogates that preserve relevant behavior, and (iii) constrains pointer-rich state spaces using semantic heap topologies. }For dynamic data structures, classic lazy initialization~\cite{visser2004test} incrementally materializes heap objects but remains largely shape-agnostic; \projName instead injects domain-informed topological partitions to reduce search space explosion while keeping data fields symbolic. 
\rev{\paragraph*{Hybrid Fuzzing and Symbolic Execution}} Compared to symbolic/concolic execution, fuzzing (e.g., AFL++~\cite{aflpp}) offers high throughput but typically struggles with deep, constraint-heavy branches~\cite{manes_fuzzing_survey}. \rev{Systems such as Driller~\cite{driller} and QSYM~\cite{qsym} are complementary to \projName{}: they combine fuzzing with concolic execution to improve input-space exploration, whereas \projName{} combines LLM guidance with symbolic execution to target solver-hostile internal fragments through selective solver-aiding ghost code rather than repeated input mutation or full-path reasoning.}
\paragraph*{LLM-based Symbolic and Concolic Execution} Recent work integrates machine learning and large language models (LLMs) into symbolic execution to mitigate scalability limits caused by complex constraints and environment modeling.
Early approaches focused on learned path exploration heuristics, using techniques such as Monte Carlo Tree Search~\cite{legion} or learned branch prioritization policies (Learch~\cite{learch}). More recent systems embed LLMs directly into symbolic or concolic execution. Solver-free approaches such as PALM~\cite{palm} and AutoBug~\cite{autobug} rely on LLMs to generate satisfying inputs for program fragments or path-wise slices, bypassing SMT solvers at the cost of global constraint consistency and scalability on deep paths. Cottontail~\cite{cottontail} targets structured-input programs using LLM-guided grammar-based input generation, while ConcoLLMic~\cite{concollmic} adopts an agentic design in which LLMs summarize traces and reason about constraints. Although LLMs improve coverage by bypassing solver bottlenecks, we showed that these systems are ineffective in handling complex program dependencies. Our hybrid approach that combines LLMs with SMT solvers addresses this problem.
\paragraph*{LLM-assisted Fuzzing and Test Generation}
LLMs have also been applied to fuzzing and test generation primarily as input generators, rather than as engines for symbolic reasoning. CodaMOSA~\cite{codamosa} uses language models to synthesize new tests when coverage plateaus, while other systems apply LLM-driven mutations to protocol and application fuzzing~\cite{chatfuzz,inputblaster}. Although effective at expanding the input space and improving shallow coverage, these approaches do not reason over execution paths or symbolic constraints and thus remain limited in reaching deep program states~\cite{manes_fuzzing_survey}.
\paragraph*{Program Inversion}
Program inversion aims to generate an inverse program that computes inputs that lead to a given output. Previous work synthesized inverse computations using symbolic execution and solvers~\cite{srivastava2011path,Abramov2000InverseComputation}, however these techniques are typically limited to small fragments or solver-friendly semantics, making them unsuitable for our use case. Just-tri-it~\cite{dai2025reducing} uses LLM-generated algorithm inverses to check cross-task solution consistency, which improves the quality of generated code. \projName applies an LLM to invert only a small program fragment, which is easier, and uses such inverses to propagate values across execution paths. \rev{Unlike prior inversion work, \projName does not require synthesizing a complete inverse program; it uses approximate fragment inverses inside a bidirectional meet-in-the-middle loop that combines LLM-generated code with SMT reasoning.}
\paragraph*{Reasoning About Programs with LLMs} Recent work has studied and improved LLMs' ability to reason about code~\cite{chen2024reasoning,liu2025evaluating}. HoarePrompt~\cite{bouras2025hoareprompt} structures chain-of-thought using ideas from strongest postcondition calculus, closely related to symbolic execution. AutoBug~\cite{autobug} improves LLM reasoning by decomposing programs into simpler components. \projName is complementary to these approaches: advances in LLM reasoning can directly improve the quality of the ghost code.

%% Recent work uses LLMs to approximate inversion when formal reasoning is
%% intractable, inferring candidate inputs through semantic reasoning
%% \cite{autobug}.
%% \projName{} adopts a solver-preserving variant of this idea, using
%% LLM-synthesized \emph{partial inverses} to propagate constraints through
%% solver-hostile fragments while validating all candidates against the original
%% program.

\section{Conclusion}

This work presents \projName, a hybrid symbolic execution framework that addresses key bottlenecks of SMT-backed symbolic execution using LLMs without giving up the ability to enforce globally consistent, cross-procedural constraints. Rather than replacing the solver with an LLM, \projName uses the LLM selectively to generate lightweight, solver-aiding ghost code tailored to solver-hostile fragments: inverse procedures for bidirectional constraint propagation, solver-friendly surrogate models that preserve control-relevant behavior, and semantic heap topologies that curb path explosion in pointer-rich code. Implemented on top of KLEE and evaluated across synthetic logic bombs, a real mathematical library, and large structured-input applications, \projName consistently improves coverage and reduces LLM cost relative to both traditional symbolic execution and prior LLM-driven alternatives.

%% \newpage
\section*{Data Availability}

All experimental artifacts required to reproduce the results of in this paper are publicly available at:
\url{https://figshare.com/s/09550df85983cbf4ae71}. 
The artifact package includes: (i) the complete implementation of \emph{\projName{}}, (ii) all baseline implementations and
  artifacts for KLEE, ConcoLLMic, and Cottontail, (iii) benchmark programs and harnesses for all six evaluated benchmark
  groups (LogicBombs, FDLibM, GNU \texttt{bc}, \texttt{jq}, \texttt{libexpat}, and \texttt{libyaml}), (iv) generated test
  cases, execution logs, and intermediate artifacts, and (v) post-processed results used to produce the tables and newly added
  analyses in the paper.

%% \section*{Acknowledgments}
%% This work was supported by the National Natural Science Foundation of China (NSFC) under Grant No.~W2542035. We thank Bo Wang for insightful discussions.




\bibliography{refs}
\bibliographystyle{plain}

\end{document}
